% Options for packages loaded elsewhere
\PassOptionsToPackage{unicode}{hyperref}
\PassOptionsToPackage{hyphens}{url}
\PassOptionsToPackage{dvipsnames,svgnames,x11names}{xcolor}
%
\documentclass[
  10pt,
]{article}
\usepackage{amsmath,amssymb}
\usepackage{iftex}
\ifPDFTeX
  \usepackage[T1]{fontenc}
  \usepackage[utf8]{inputenc}
  \usepackage{textcomp} % provide euro and other symbols
\else % if luatex or xetex
  \usepackage{unicode-math} % this also loads fontspec
  \defaultfontfeatures{Scale=MatchLowercase}
  \defaultfontfeatures[\rmfamily]{Ligatures=TeX,Scale=1}
\fi
\usepackage{lmodern}
\ifPDFTeX\else
  % xetex/luatex font selection
  \setmainfont[]{DejaVu Serif}
  \setsansfont[]{DejaVu Sans}
  \setmonofont[]{DejaVu Sans Mono}
  \setmathfont[]{Latin Modern Math}
\fi
% Use upquote if available, for straight quotes in verbatim environments
\IfFileExists{upquote.sty}{\usepackage{upquote}}{}
\IfFileExists{microtype.sty}{% use microtype if available
  \usepackage[]{microtype}
  \UseMicrotypeSet[protrusion]{basicmath} % disable protrusion for tt fonts
}{}
\makeatletter
\@ifundefined{KOMAClassName}{% if non-KOMA class
  \IfFileExists{parskip.sty}{%
    \usepackage{parskip}
  }{% else
    \setlength{\parindent}{0pt}
    \setlength{\parskip}{6pt plus 2pt minus 1pt}}
}{% if KOMA class
  \KOMAoptions{parskip=half}}
\makeatother
\usepackage{xcolor}
\usepackage[margin=2.0cm]{geometry}
\usepackage{color}
\usepackage{fancyvrb}
\newcommand{\VerbBar}{|}
\newcommand{\VERB}{\Verb[commandchars=\\\{\}]}
\DefineVerbatimEnvironment{Highlighting}{Verbatim}{commandchars=\\\{\}}
% Add ',fontsize=\small' for more characters per line
\newenvironment{Shaded}{}{}
\newcommand{\AlertTok}[1]{\textcolor[rgb]{1.00,0.00,0.00}{\textbf{#1}}}
\newcommand{\AnnotationTok}[1]{\textcolor[rgb]{0.38,0.63,0.69}{\textbf{\textit{#1}}}}
\newcommand{\AttributeTok}[1]{\textcolor[rgb]{0.49,0.56,0.16}{#1}}
\newcommand{\BaseNTok}[1]{\textcolor[rgb]{0.25,0.63,0.44}{#1}}
\newcommand{\BuiltInTok}[1]{\textcolor[rgb]{0.00,0.50,0.00}{#1}}
\newcommand{\CharTok}[1]{\textcolor[rgb]{0.25,0.44,0.63}{#1}}
\newcommand{\CommentTok}[1]{\textcolor[rgb]{0.38,0.63,0.69}{\textit{#1}}}
\newcommand{\CommentVarTok}[1]{\textcolor[rgb]{0.38,0.63,0.69}{\textbf{\textit{#1}}}}
\newcommand{\ConstantTok}[1]{\textcolor[rgb]{0.53,0.00,0.00}{#1}}
\newcommand{\ControlFlowTok}[1]{\textcolor[rgb]{0.00,0.44,0.13}{\textbf{#1}}}
\newcommand{\DataTypeTok}[1]{\textcolor[rgb]{0.56,0.13,0.00}{#1}}
\newcommand{\DecValTok}[1]{\textcolor[rgb]{0.25,0.63,0.44}{#1}}
\newcommand{\DocumentationTok}[1]{\textcolor[rgb]{0.73,0.13,0.13}{\textit{#1}}}
\newcommand{\ErrorTok}[1]{\textcolor[rgb]{1.00,0.00,0.00}{\textbf{#1}}}
\newcommand{\ExtensionTok}[1]{#1}
\newcommand{\FloatTok}[1]{\textcolor[rgb]{0.25,0.63,0.44}{#1}}
\newcommand{\FunctionTok}[1]{\textcolor[rgb]{0.02,0.16,0.49}{#1}}
\newcommand{\ImportTok}[1]{\textcolor[rgb]{0.00,0.50,0.00}{\textbf{#1}}}
\newcommand{\InformationTok}[1]{\textcolor[rgb]{0.38,0.63,0.69}{\textbf{\textit{#1}}}}
\newcommand{\KeywordTok}[1]{\textcolor[rgb]{0.00,0.44,0.13}{\textbf{#1}}}
\newcommand{\NormalTok}[1]{#1}
\newcommand{\OperatorTok}[1]{\textcolor[rgb]{0.40,0.40,0.40}{#1}}
\newcommand{\OtherTok}[1]{\textcolor[rgb]{0.00,0.44,0.13}{#1}}
\newcommand{\PreprocessorTok}[1]{\textcolor[rgb]{0.74,0.48,0.00}{#1}}
\newcommand{\RegionMarkerTok}[1]{#1}
\newcommand{\SpecialCharTok}[1]{\textcolor[rgb]{0.25,0.44,0.63}{#1}}
\newcommand{\SpecialStringTok}[1]{\textcolor[rgb]{0.73,0.40,0.53}{#1}}
\newcommand{\StringTok}[1]{\textcolor[rgb]{0.25,0.44,0.63}{#1}}
\newcommand{\VariableTok}[1]{\textcolor[rgb]{0.10,0.09,0.49}{#1}}
\newcommand{\VerbatimStringTok}[1]{\textcolor[rgb]{0.25,0.44,0.63}{#1}}
\newcommand{\WarningTok}[1]{\textcolor[rgb]{0.38,0.63,0.69}{\textbf{\textit{#1}}}}
\usepackage{longtable,booktabs,array}
\usepackage{calc} % for calculating minipage widths
% Correct order of tables after \paragraph or \subparagraph
\usepackage{etoolbox}
\makeatletter
\patchcmd\longtable{\par}{\if@noskipsec\mbox{}\fi\par}{}{}
\makeatother
% Allow footnotes in longtable head/foot
\IfFileExists{footnotehyper.sty}{\usepackage{footnotehyper}}{\usepackage{footnote}}
\makesavenoteenv{longtable}
\setlength{\emergencystretch}{3em} % prevent overfull lines
\providecommand{\tightlist}{%
  \setlength{\itemsep}{0pt}\setlength{\parskip}{0pt}}
\setcounter{secnumdepth}{-\maxdimen} % remove section numbering
% doc_header.tex -- LaTeX preamble additions used by build_docs.sh for every
% Markdown -> LaTeX -> PDF conversion (supplements, README, submission files).
\usepackage{etoolbox}
% longtable may not be loaded by pandoc when a document has no table
\usepackage{longtable}
% bold vectors: \bm -> \symbf under unicode-math (XeLaTeX), bm package otherwise
\makeatletter
\@ifpackageloaded{unicode-math}{\providecommand{\bm}[1]{\symbf{#1}}}{\usepackage{bm}}
\makeatother
% smaller type in long tables so that wide audit tables fit the text width;
% inside a table a line break is allowed after an underscore of plain text
% (pandoc writes it as \_), so that long identifiers wrap inside their column
% rows are set a little apart, so that symbols with indices do not touch
\AtBeginEnvironment{longtable}{\footnotesize\renewcommand{\arraystretch}{1.45}%
  \let\sptunderscore\_\def\_{\sptunderscore\allowbreak}}
\setlength{\LTcapwidth}{\textwidth}
% allow line breaks after slashes and hyphens inside table cells
\usepackage{xurl}
% code blocks: break over-long lines instead of running into the margin.
% Plain ``` blocks are typeset by pandoc as {verbatim}, highlighted ones
% (```bash) as {Highlighting}; both are re-declared as breakable fancyvrb
% environments.  The continuation line is marked with a small hooked arrow.
\usepackage{fvextra}
\RecustomVerbatimEnvironment{verbatim}{Verbatim}{breaklines,breakanywhere,fontsize=\small}
\makeatletter
\@ifundefined{Highlighting}{}{%
  \RecustomVerbatimEnvironment{Highlighting}{Verbatim}{breaklines,breakanywhere,fontsize=\small,commandchars=\\\{\}}}
\makeatother
% inline code (pandoc writes it as \texttt{...}): never hyphenate -- an inserted
% hyphen could be read as part of a file name -- but allow a line break after
% an underscore, slash, comma, colon or equals sign, so that long paths and
% column lists wrap inside the text width and inside table columns.
% Inside a formula (e.g. the name of a notebook symbol in an equation) the same
% command must first leave math mode, which \text of amsmath does.
\ExplSyntaxOn
\tl_new:N \l__spt_code_tl
\cs_new_protected:Npn \__spt_texttt:n #1
  {
    \tl_set:Nn \l__spt_code_tl {#1}
    \regex_replace_all:nnN { ([\/\,\:\=]) } { \1 \c{allowbreak} } \l__spt_code_tl
    \regex_replace_all:nnN { \c{_} } { \c{_} \c{allowbreak} } \l__spt_code_tl
    \group_begin: \ttfamily \hyphenchar \font = -1 \scan_stop: \l__spt_code_tl \group_end:
  }
\RenewDocumentCommand \texttt { m }
  {
    \mode_if_math:TF { \text { \__spt_texttt:n {#1} } } { \__spt_texttt:n {#1} }
  }
\ExplSyntaxOff
% A figure is written in the Markdown source as an image paragraph followed by
% a caption paragraph ("**Figure S1.** ..."), so that the caption is also
% visible on GitHub.  An image that starts a paragraph is centered.
\usepackage{graphicx}% pandoc loads it only for documents that contain an image
\NewCommandCopy\sptincludegraphics\includegraphics
\RenewDocumentCommand\includegraphics{O{} m}{%
  \ifvmode\begin{center}\sptincludegraphics[#1]{#2}\end{center}%
  \else\sptincludegraphics[#1]{#2}\fi}
% keep figures where they are referenced
\usepackage{float}
\floatplacement{figure}{H}
% running header with the document title
\usepackage{fancyhdr}
\pagestyle{fancy}
\fancyhf{}
\fancyhead[L]{\small PROJECT spacetimes, Revision\_Deflation\_CJT\_and\_Kohn-Sham}
\fancyhead[R]{\small\thepage}
\renewcommand{\headrulewidth}{0.4pt}
% 16x16 matrices: tighter columns so that a full matrix fits the text width
\setlength{\arraycolsep}{3.2pt}
% The documents of this project contain long tagged symbols, file paths and
% identifiers that cannot be hyphenated.  A paragraph is therefore allowed to
% have looser lines before TeX gives up and lets a line stick out of the margin.
\AtBeginDocument{\tolerance=2000 \emergencystretch=5em \hbadness=2000}
\ifLuaTeX
  \usepackage{selnolig}  % disable illegal ligatures
\fi
\IfFileExists{bookmark.sty}{\usepackage{bookmark}}{\usepackage{hyperref}}
\IfFileExists{xurl.sty}{\usepackage{xurl}}{} % add URL line breaks if available
\urlstyle{same}
\hypersetup{
  pdftitle={MASTER provenance of the symbols of PROJECT spacetimes},
  pdfauthor={Patrick Lee Nash, Ph.D.},
  colorlinks=true,
  linkcolor={blue},
  filecolor={Maroon},
  citecolor={blue},
  urlcolor={blue},
  pdfcreator={LaTeX via pandoc}}

\title{MASTER provenance of the symbols of PROJECT spacetimes}
\usepackage{etoolbox}
\makeatletter
\providecommand{\subtitle}[1]{% add subtitle to \maketitle
  \apptocmd{\@title}{\par {\large #1 \par}}{}{}
}
\makeatother
\subtitle{Every symbol of the project with one meaning, and the proof
that no symbol has two}
\author{Patrick Lee Nash, Ph.D.}
\date{November 13, 2026}

\begin{document}
\maketitle

\hypertarget{what-this-document-is}{%
\section{1. What this document is}\label{what-this-document-is}}

The notebooks, papers and documents on which PROJECT spacetimes rests
were written over many years, and they often use the same letter for
different objects. The letter \(g\) alone stands, in them, for a
component of a metric, for a rate of deflation, for a coupling of a
self-interaction and for the coupling of a scalar field; the letter
\(M\) for a mass and for a family of matrices; the letter \(F\) for a
matrix of numbers and for scattering amplitudes. A reader who meets such
a letter cannot know what it means without knowing where the sentence
came from.

This document removes that ambiguity for the whole project. It contains

\begin{enumerate}
\def\labelenumi{\arabic{enumi}.}
\tightlist
\item
  the \textbf{rules} by which symbols of the project are told apart
  (Section 2);
\item
  the list of the \textbf{collisions} that were found, with the decision
  taken for each (Section 3);
\item
  a \textbf{dictionary} from the notation of the author's document on
  Kohn--Sham matter and gravity on the 7+1 spacetime to the symbols of
  the project (Section 4);
\item
  the \textbf{proof} that the registry is free of collisions: seven
  machine checks, with their complete output (Section 5), and the
  instructions for repeating them (Section 6);
\item
  the \textbf{complete registry}: 229 entries in 24 groups, with 339
  names used in programs and data files (Section 7), and the key to the
  column ``origin'' (Section 8).
\end{enumerate}

Every symbol that occurs in this document is defined in Section 7 of
this document; nothing has to be looked up elsewhere.

\textbf{How the registry is used.} The registry is one table,
\texttt{provenance/symbols\_master.csv}. A script writes it; a second
script checks it and checks the documents and the data files of the
folder against it. Every document of the folder
\texttt{Revision\_Deflation\_CJT\_and\_Kohn-Sham} uses the symbols of
the registry with the meanings given there, and a symbol that is not in
the registry may not be introduced without being added to it.

\hypertarget{the-rules}{%
\section{2. The rules}\label{the-rules}}

\textbf{R1. A tag says which spacetime an object belongs to.} Objects of
the flat 4+4 spacetime carry the tag 1644 (\(\gamma_{1644}^A\),
\(C^-_{1644}\), \(W_{1644}\)), objects of the flat 7+1 spacetime the tag
1671 (\(\gamma_{1671}^A\), \(\chi_{1671}\)). The fields on the deflating
spacetimes carry the tags 1644d and 1671d. Where a section treats one
spacetime only, the tag may be dropped; the section says so.

\textbf{R2. A suffix says which sector an object belongs to.} Every
quantity of the Vorton-Scalar carries the suffix VS: \(\Phi_{VS}\),
\(v_{VS}\), \(y_{VS}\), \(\lambda_{VS}\), \(g_{VS}\), \(\alpha_{VS}\),
\(D_{VS}\), \(G_{VS}\).

\textbf{R3. The four kinds of \(g\) are written in four different ways.}

\begin{longtable}[]{@{}
  >{\raggedright\arraybackslash}p{(\columnwidth - 4\tabcolsep) * \real{0.3333}}
  >{\raggedright\arraybackslash}p{(\columnwidth - 4\tabcolsep) * \real{0.3333}}
  >{\raggedright\arraybackslash}p{(\columnwidth - 4\tabcolsep) * \real{0.3333}}@{}}
\toprule\noalign{}
\begin{minipage}[b]{\linewidth}\raggedright
object
\end{minipage} & \begin{minipage}[b]{\linewidth}\raggedright
how it is written
\end{minipage} & \begin{minipage}[b]{\linewidth}\raggedright
example
\end{minipage} \\
\midrule\noalign{}
\endhead
\bottomrule\noalign{}
\endlastfoot
rate of deflation (the author's names) & upright, two digits, never a
subscript & \(\mathrm{g55}\), \(\mathrm{g66}\), \(\mathrm{g77}\) \\
component of the metric of a deflating spacetime & italic, superscript
\((4+4)\) or \((7+1)\), two indices & \(g^{(4+4)}_{55}\),
\(g^{(7+1)}_{\mu\nu}\) \\
coupling of the self-interaction & italic, one digit or the letter \(r\)
as subscript & \(g_0,\dots,g_8\), \(g_r\), \(g_0^{\rm eff}\) \\
coupling of the Vorton-Scalar & italic, suffix VS & \(g_{VS}\) \\
\end{longtable}

\textbf{R4. Type distinguishes matrices of numbers from amplitudes, and
scale factors from indices.} Sans-serif type: the Fierz matrix
\(\mathsf F\), the density matrix \(\mathsf n\), the Hamiltonian matrix
\(\mathsf h\), the two scale factors \(\mathsf a\), \(\mathsf b\) of the
equations of gravity. Italic type: the amplitudes \(F_S\), \(F_P\),
\(F_0\), \(F_1\), \(F_7\), \(F_8\); the indices \(a\), \(b\). Script
type: the Einstein tensor \(\mathcal G\), the Lanczos tensor
\(\mathcal H\), the Killing vectors \(\mathcal X\), the sources
\(\mathcal D,\mathcal T,\mathcal P,\mathcal A\) of the perturbation
equations.

\textbf{R5. An index, an argument or a bar distinguishes objects that
share a letter.} \(\Lambda\) without index is a momentum cutoff,
\(\Lambda_8\) and \(\bar\Lambda\) are cosmological terms; \(\kappa_8\)
is the gravitational coupling and \(\kappa(k)\) a pair correlation;
\(E_F\), \(E_k\), \(E_\perp\), \(E^*\) are energies and
\(E^{(k)\mu}{}_\nu\), with the order in parentheses, is a Lovelock
tensor; \(\alpha^j\) with a superscript are Dirac matrices,
\(\alpha_1,\alpha_2,\alpha_3\) with a digit as subscript are couplings
of the Lovelock terms, \(\alpha_{VS}\) and \(\alpha_c\) are coupling
strengths; \(\sigma_8\) is the amplitude of clustering, \(\sigma_v\) a
velocity dispersion and \(\sigma(p)\) a self-energy.

\textbf{R6. A letter that the author uses keeps the author's meaning},
wherever that is possible without a collision: \(G_S\), \(G_V\),
\(V_0\), \(I_r\), \(\Theta^K\), \(g_r\) of his document on Kohn--Sham
matter are used here for the same objects. Where a letter of his
collides with a symbol of the project, the registry gives the object a
new symbol and records his letter (Section 4).

\textbf{R7. Program names.} A name used in a program or as a column of a
data file is registered with its symbol. A name may belong to one entry
only. Names of columns that are plain words (``content'', ``case'') are
listed in the checking script.

\hypertarget{the-collisions-that-were-found-and-how-each-was-resolved}{%
\section{3. The collisions that were found, and how each was
resolved}\label{the-collisions-that-were-found-and-how-each-was-resolved}}

\begin{longtable}[]{@{}
  >{\raggedright\arraybackslash}p{(\columnwidth - 4\tabcolsep) * \real{0.3333}}
  >{\raggedright\arraybackslash}p{(\columnwidth - 4\tabcolsep) * \real{0.3333}}
  >{\raggedright\arraybackslash}p{(\columnwidth - 4\tabcolsep) * \real{0.3333}}@{}}
\toprule\noalign{}
\begin{minipage}[b]{\linewidth}\raggedright
letter
\end{minipage} & \begin{minipage}[b]{\linewidth}\raggedright
the different objects that used it
\end{minipage} & \begin{minipage}[b]{\linewidth}\raggedright
resolution in the registry
\end{minipage} \\
\midrule\noalign{}
\endhead
\bottomrule\noalign{}
\endlastfoot
\(g\) & rate of deflation; component of the metric; coupling of the
self-interaction; coupling of the Vorton-Scalar & rule R3 \\
\(g_{1644d}\), \(g_{1671d}\) & the metric of the first revision on the
deflating spacetimes, in which the time-like component was positive on
7+1 & retired; the metric is \(g^{(4+4)}\), \(g^{(7+1)}\) in the
author's sign convention (space-like components positive) \\
\(\eta\) & flat metric; matrix of the squares of the Dirac matrices;
function \(\eta\) of the synchronous gauge; sign \(\eta_K\) of the
author's document; loading of an evolution & \(\eta^{(4+4)}\),
\(\eta^{(7+1)}\) (metric); \(\eta_{1644}\), \(\eta_{1671}\) (squares of
the matrices); \(\eta_{\rm sy}\); \(\mathrm{sgn}_K\); the word
``loading'' \\
\(F\) & Fierz matrix \(F_{rs}\); pairing amplitudes \(F_S\), \(F_P\),
\(F_0\) & \(\mathsf F_{rr'}\) in sans-serif type for the matrix \\
\(\sigma\) & exponent \(\sigma_r\) of the phase of the author's
document; amplitude \(\sigma_8\); velocity dispersion; self-energy
\(\sigma(p)\) & phase \(\zeta_r\); \(\sigma_8\); \(\sigma_v\);
\(\sigma(p)\) \\
\(M\) & mass of the particles in matter; vertex matrices \(M_K\); gap
matrix & mass \(m^*\); the vertex matrices are written out as
\(\beta\Theta^K\); gap matrix \(\mathsf M\) \\
\(\rho\) & number density of matter; energy density of cosmology;
coordinate index & number density \(n\) and scalar density \(n_S\);
\(\rho\) is a mass or energy density of cosmology; as an index only as
sub- or superscript \\
\(\Pi\) & interaction part of the pressure & \(p_{\rm int}\) \\
\(\kappa\) & gravitational coupling; pair correlation; ratios
\(\kappa_S\), \(\kappa_V\) of couplings; bound on a cross-section &
\(\kappa_8\); \(\kappa(k)\); the ratios are written \(-G_S/C\),
\(G_V/C\); \(\kappa_{\rm SI}\) \\
\(\Lambda\) & momentum cutoff; range of the contact couplings;
cosmological terms & \(\Lambda\) (cutoff); range parameter
\(m_{\rm ct}\); \(\Lambda_8\), \(\bar\Lambda\) \\
\(P\) & polarization of the two halves of the 4+4 field; generator of
\(U(1)_E\); chiral projectors; momenta of a classical particle &
\(\mathrm{pol}_W\); \(P\), \(P_c\) (generators); \(P^{L,R}\);
\(P_\mu\) \\
\(a\), \(b\) & indices of the extra axes; scale factors of the equations
of gravity; scale factor of ordinary space & indices \(a,b\);
\(\mathsf a\), \(\mathsf b\); \(a_{\rm sp}\) \\
\(s_t\) & stretch factors of the momenta along the deflating coordinates
& written as the ratio \(b_t^{\rm ref}/b_t\) \\
\(\tau\) & deflation clock; clock of the undisturbed deflation;
conformal time; optical depth & \(\tau_d\); \(\tau_d^0\);
\(\tau_{\rm cf}\); \(\tau\) (optical depth, only in the list of
cosmological parameters) \\
\(K\) & generator of the evolution on 4+4; index list of the
self-interaction; sectional curvature & \(K\), \(K_\perp\) (no index or
the index \(\perp\)); \(K\), \(J\) as index lists appear only as sub- or
superscripts of \(\Theta\) and under sums; \(K_{AB}\) with two
indices \\
\(x\) & coordinates \(x0,\dots,x7\); ratio \(m^*/m\); ratio
\(\phi/v_{VS}\) & coordinates always with their digit; \(x=\phi/v_{VS}\)
(equal to \(m^*/m\) only without bare mass and contact couplings) \\
\(\Phi\) & Vorton-Scalar; state of a pair & \(\Phi_{VS}\); \(\Phi\),
\(\Phi(\vec p)\) without suffix \\
\(\Gamma\) & Christoffel symbols; Dirac matrices of the author's 7+1
document; effective action; Euler's gamma function &
\(\Gamma^\rho{}_{\mu\nu}\) (three indices); Dirac matrices are
\(\gamma\); \(\Gamma[\ldots]\) with its argument in brackets;
\(\Gamma_{\rm E}\) \\
\(\omega\) & product of the eight Dirac matrices; coefficients of the
spin connection; energy of a particle in the author's document;
densities of cosmology & \(\omega_{1644}\), \(\omega_{1671}\);
\(\omega_{\mu AB}\) (three indices); \(E^*(k)\); \(\omega_b\),
\(\omega_c\) \\
\end{longtable}

\hypertarget{dictionary-the-authors-document-on-kohnsham-matter-and-gravity-on-71}{%
\section{4. Dictionary: the author's document on Kohn--Sham matter and
gravity on
7+1}\label{dictionary-the-authors-document-on-kohnsham-matter-and-gravity-on-71}}

The author's document that defines, solves and analyzes the 7+1 problem
with self-interaction uses its own conventions. The calculations of this
project reproduce its numbers; the table says how its notation is
translated.

\begin{longtable}[]{@{}
  >{\raggedright\arraybackslash}p{(\columnwidth - 4\tabcolsep) * \real{0.3333}}
  >{\raggedright\arraybackslash}p{(\columnwidth - 4\tabcolsep) * \real{0.3333}}
  >{\raggedright\arraybackslash}p{(\columnwidth - 4\tabcolsep) * \real{0.3333}}@{}}
\toprule\noalign{}
\begin{minipage}[b]{\linewidth}\raggedright
in the author's document
\end{minipage} & \begin{minipage}[b]{\linewidth}\raggedright
in this project
\end{minipage} & \begin{minipage}[b]{\linewidth}\raggedright
remark
\end{minipage} \\
\midrule\noalign{}
\endhead
\bottomrule\noalign{}
\endlastfoot
time coordinate \(x_0\); ordinary space \(x_1,x_2,x_3\); auxiliary space
\(x_4,\dots,x_7\) & time \(x4\); ordinary space \(x1,x2,x3\); the four
other space-like coordinates \(x0,x5,x6,x7\) & the exchange of the index
values 0 and 4 (\(\mathrm{swap}_{04}\)) \\
\(\eta=\mathrm{diag}(+1,-1,\dots,-1)\): squares of the Dirac matrices
and signs of the metric & \(\eta_{1671}\) with the time at index 4; the
metric of the deflating spacetime is \(g^{(7+1)}\) with space-like
components positive & \(\eta^{(7+1)}=-\eta_{1671}\) \\
Dirac matrices \(\Gamma_A\) & \(\gamma_{1671}^A\) & the same real
matrices \\
Christoffel symbols in braces & \(\Gamma^\rho{}_{\mu\nu}\) & \\
couplings \(g_r\); invariants \(I_r\); matrices \(\Theta^K\),
\(\Theta_K\) & the same symbols & \\
sign \(\eta_K\); phase \(i^{\sigma_r}\) & \(\mathrm{sgn}_K\);
\(\zeta_r\) & \\
Fierz matrix \(F_{rs}\) & \(\mathsf F_{rr'}\) & \\
density matrix \(R\) & \(\mathsf n\,\beta\) & \\
Hartree--Fock couplings \(G_0\), \(G_1\) & \(G_S^{\rm HF}\),
\(G_V^{\rm HF}\) & \\
\(G_S\), \(G_V\), \(V_0\) & the same symbols & the mean-field couplings
and the vector field \\
mass in matter \(M\) & \(m^*\) & the data files keep the name
\texttt{M} \\
densities \(\rho\), \(\rho_S\) & \(n\), \(n_S\) & the data files keep
\texttt{rho}, \texttt{rho\_S} \\
energy of a particle \(\omega\) & \(E^*(k)\) & \\
interaction pressure \(\Pi\) & \(p_{\rm int}\) & the data files keep
\texttt{Pi} \\
scale factors \(a\), \(b\) and rates \(H_a\), \(H_b\) & \(\mathsf a\),
\(\mathsf b\), \(H_{\mathsf a}\), \(H_{\mathsf b}\) & the data files
keep \texttt{a}, \texttt{b}, \texttt{H\_a}, \texttt{H\_b} \\
gravitational coupling \(\kappa\) & \(\kappa_8\) & \\
Lovelock tensors \(E^{(k)}\), couplings \(\alpha_k\), cosmological term
\(\bar\Lambda\) & the same symbols & \\
\end{longtable}

\hypertarget{the-proof-seven-machine-checks}{%
\section{5. The proof: seven machine
checks}\label{the-proof-seven-machine-checks}}

The script \texttt{code/scripts/check\_symbols.py} reads the registry,
every document of the folder and every data file, and verifies seven
statements. Its complete output:

\begin{verbatim}
PROJECT spacetimes, Revision_Deflation_CJT_and_Kohn-Sham: checks of the symbols (MASTER registry provenance/symbols_master.csv)
generated by: code/scripts/check_symbols.py
7 of 7 checks passed

[PASS] MS1  the 229 symbol texts of the MASTER registry are pairwise different, and so are its 339 names in programs and data files
[PASS] MS2  no definition of the MASTER registry sends the reader to another document (searched for the names of the documents of the project and for the words "main registry", "supplement", "see the")
[PASS] MS3  none of the 41 names in programs that this sub-project introduces itself (entries of the groups 10 to 15 with the origin "new"; names of more than two characters) is a symbol of one of the author's notebooks (1921 symbols compared). Names that are the author's own (G_S, G_V, V_0, p_a, p_b and others of his Kohn-Sham sections) are used for the objects for which he uses them
[PASS] MS4  in the 35 documents of this folder the rates g55, g66, g77 are written upright 556 times and never as a letter with a subscript; each of the 72 components of the metric g_{55}, g_{66}, g_{77}, g_{tt} carries the superscript (4+4) or (7+1) or stands in a paragraph that names it a component of the metric
[PASS] MS5  the retired symbols do not occur in the documents: \varkappa nowhere; g_{1644d}, g_{1671d} only where a paragraph says that they belong to the earlier sign convention
[PASS] MS6  the 14 retired names of programs (tau, kappa_t, Q_over_eps, Omega0, S2, mu and nu as orders of Bessel functions, A, Xi without index, m_b, w0, omega_c for the shell, the metric names g_1644d and g_1671d of the earlier sign convention, and Lambda as the range of the contact couplings) do not occur in the 132 data files that the routines of this folder write
[PASS] MS7  each of the 1093 columns of the 111 data files of this folder is a registered name, a column of the flat-space part of the project, one of the 249 descriptive words listed in the script, or (213 columns) a systematic name
\end{verbatim}

What the checks mean:

\begin{itemize}
\tightlist
\item
  \textbf{MS1} No two entries of the registry have the same symbol text,
  and no name of a program or data file belongs to two entries. This is
  the statement that no symbol has two meanings.
\item
  \textbf{MS2} No definition sends the reader to another document.
\item
  \textbf{MS3} No name that this sub-project introduces is already a
  symbol of one of the author's notebooks with another meaning.
\item
  \textbf{MS4} In the documents, a rate is never written as a letter
  with a subscript, and a component of the metric always carries its
  superscript.
\item
  \textbf{MS5}, \textbf{MS6} The retired symbols and names do not occur.
\item
  \textbf{MS7} Every column of every data file is a registered name.
\end{itemize}

\hypertarget{how-to-repeat-the-checks}{%
\section{6. How to repeat the checks}\label{how-to-repeat-the-checks}}

Run from the folder \texttt{Revision\_Deflation\_CJT\_and\_Kohn-Sham} of
the repository:

\begin{Shaded}
\begin{Highlighting}[]
\ExtensionTok{python3}\NormalTok{ code/scripts/write\_symbol\_registry.py}
\ExtensionTok{python3}\NormalTok{ code/scripts/check\_symbols.py}
\end{Highlighting}
\end{Shaded}

The first command writes \texttt{provenance/symbols\_master.csv} from
three sources: the registry of the flat-spacetime part of the project
(\texttt{../provenance/symbols.csv}), the registry of the first revision
on the deflating spacetimes
(\texttt{../Revision\_Deflation/provenance/symbols\_deflation.csv}), and
the list of new and reworded entries inside the script. It stops with an
error if two entries would have the same symbol or the same program
name. The second command performs the seven checks and writes
\texttt{data/symbol\_checks.txt}; its exit status is zero only if all
pass.

\hypertarget{the-complete-registry}{%
\section{7. The complete registry}\label{the-complete-registry}}

Columns: the symbol as it is written in the documents; the names that
programs and data files use for it; its meaning; its origin (the key is
in Section 8; ``new'' marks a symbol introduced by this project,
``author'' a name fixed by an instruction of the author).

\textbf{1 Coordinates and indices}

\begin{longtable}[]{@{}
  >{\raggedright\arraybackslash}p{(\columnwidth - 6\tabcolsep) * \real{0.1605}}
  >{\raggedright\arraybackslash}p{(\columnwidth - 6\tabcolsep) * \real{0.1235}}
  >{\raggedright\arraybackslash}p{(\columnwidth - 6\tabcolsep) * \real{0.5309}}
  >{\raggedright\arraybackslash}p{(\columnwidth - 6\tabcolsep) * \real{0.1852}}@{}}
\toprule\noalign{}
\begin{minipage}[b]{\linewidth}\raggedright
symbol
\end{minipage} & \begin{minipage}[b]{\linewidth}\raggedright
names in programs and data files
\end{minipage} & \begin{minipage}[b]{\linewidth}\raggedright
meaning
\end{minipage} & \begin{minipage}[b]{\linewidth}\raggedright
origin
\end{minipage} \\
\midrule\noalign{}
\endhead
\bottomrule\noalign{}
\endlastfoot
\(X=\{x0,\dots,x7\}\) & \texttt{x0} \texttt{x1} \texttt{x2} \texttt{x3}
\texttt{x4} \texttt{x5} \texttt{x6} \texttt{x7} & the eight coordinates
of both spacetimes, never renamed: \(x0\) hidden space-like;
\(x1,x2,x3\) ordinary space; \(x4\) ordinary time; \(x5,x6,x7\)
additional coordinates (time-like in 4+4, space-like in 7+1) & N01 (text
at its head) and N04 \texttt{X} (section DEFCHART): the same eight names
with the same meaning. N02 uses the same eight names with \texttt{x0} as
the time: related by swap04 \\
\(\mathrm{swap}_{04}\) & \texttt{swap04} & the exchange of the index
values 0 and 4; it carries the labels of \texttt{7+1.nb} (time at index
0) into those of the project (time at index 4) & new \\
\(A,B,C,D\) & --- & index with the values \(0,\dots,7\); always labels
the coordinate \(xA\) & new \\
\(a,b\) & --- & indices whose range is stated at each place of use: the
four extra axes, values 0, 5, 6, 7 (in everything that concerns pairs
and the order parameter), or the four flat coordinates other than the
time, values 0, 1, 2, 3 (in the statements of the machine checks). The
letters \(\mathsf a\), \(\mathsf b\) in sans-serif type are the scale
factors of the Einstein-Lovelock evolution; \(a_{\rm sp}\),
\(a_{\rm sc}\), \(a(T)\), \(b_t\) and \(b\) have their own entries &
new \\
\(i,j\) & --- & index of space-like coordinates; values 1, 2, 3
(ordinary space) unless the range is stated: \(j=0,1,2,3\) for the four
flat coordinates other than the time where the deflating coordinates
occur as well, and all seven space-like directions in the Hamiltonian of
7+1. In the Einstein-Lovelock equations \(i,j\) label the seven space
directions & new \\
\(h\) (index) & --- & index with the values 1, 2, 3 that labels the
\(4\times4\) matrices \(s_h\), \(t_h\) & N01, N02, N05 \texttt{h} in
\texttt{s4by4{[}h{]}}, \texttt{t4by4{[}h{]}} \\
\end{longtable}

\textbf{2 Matrices of the 4+4 spacetime (tag 1644)}

\begin{longtable}[]{@{}
  >{\raggedright\arraybackslash}p{(\columnwidth - 6\tabcolsep) * \real{0.1605}}
  >{\raggedright\arraybackslash}p{(\columnwidth - 6\tabcolsep) * \real{0.1235}}
  >{\raggedright\arraybackslash}p{(\columnwidth - 6\tabcolsep) * \real{0.5309}}
  >{\raggedright\arraybackslash}p{(\columnwidth - 6\tabcolsep) * \real{0.1852}}@{}}
\toprule\noalign{}
\begin{minipage}[b]{\linewidth}\raggedright
symbol
\end{minipage} & \begin{minipage}[b]{\linewidth}\raggedright
names in programs and data files
\end{minipage} & \begin{minipage}[b]{\linewidth}\raggedright
meaning
\end{minipage} & \begin{minipage}[b]{\linewidth}\raggedright
origin
\end{minipage} \\
\midrule\noalign{}
\endhead
\bottomrule\noalign{}
\endlastfoot
\(\eta_{1644}\) & \texttt{eta1644} & metric of the flat 4+4 spacetime
and matrix of the squares of its Dirac matrices,
\(\mathrm{diag}(+1,+1,+1,+1,-1,-1,-1,-1)\): space-like \(=+1\). It
equals the flat metric \(\eta^{(4+4)}\) of the author's sign convention
for the deflating spacetimes & N01, N02, N05, N10--N12 \texttt{η4488};
in N01 also \texttt{η8} and \texttt{g\_AB}; D08 \(G_{AB}\) with the
index value 8 in the place of the value 0 \\
\(s_h\), \(t_h\) & --- & the three self-dual and the three
anti-self-dual real antisymmetric \(4\times4\) matrices & N01, N02, N05,
N10--N12 \texttt{s4by4{[}h{]}}, \texttt{t4by4{[}h{]}}; the letters are
those of D06 (\(s^h\) self-dual, \(t^h\) anti-self-dual) \\
\(\tau_{1644}^A\), \(\bar\tau_{1644}^A\) & \texttt{tau1644}
\texttt{taubar1644} & the real \(8\times8\) blocks of the 4+4 Dirac
matrices (reduced generators of the split octonion algebra) & N01,
N10--N12 \texttt{τ{[}A{]}} and \texttt{τ} with an overbar; N02
\texttt{τ{[}A{]}}, \texttt{τbar{[}A{]}}; matrices of this kind are the
\(\tau^A\), \(\bar\tau^A\) of D07 and D08, whose index runs from 1 to 8
(the value 8 there is the value 0 here) \\
\(\sigma_{88}\) & \texttt{sigma88} & the \(8\times8\) matrix with the
\(4\times4\) blocks \((0,1_4;\,1_4,0)\) & N01 \texttt{σ}, \texttt{σ88};
N02 \texttt{σ}; N05 \texttt{σ88} \\
\(\gamma_{1644}^A\) & \texttt{gamma1644} & the eight real \(16\times16\)
Dirac matrices of the 4+4 spacetime & N01, N10--N12
\texttt{T16\^{}A{[}A{]}}; N02 \texttt{T16{[}A{]}}; the \(t^A\) of D08
are built in the same way \\
\(\Sigma_{1644}^{AB}\) & \texttt{Sigma1644} &
\(\tfrac14[\gamma_{1644}^A,\gamma_{1644}^B]\), the scaled commutators;
they span \(\mathfrak{so}(4,4)\) & new name; the same matrices are
\texttt{SAB} and \texttt{sAB{[}A{]}{[}B{]}} of N10--N12; NOT
\texttt{sAB} of N02, which is a sign \\
\(\omega_{1644}\) & \texttt{omega1644} & the product
\(\gamma_{1644}^0\cdots\gamma_{1644}^7\); real chirality matrix of 4+4,
\(\mathrm{diag}(-1_8,+1_8)\) & N01 \texttt{T16\^{}A{[}8{]}}; N02
\texttt{T16{[}8{]}} \\
\(P_{1644}^{L}\), \(P_{1644}^{R}\) & \texttt{PL1644} \texttt{PR1644} &
the real chiral projectors \(\tfrac12(1\mp\omega_{1644})\) & N01
\texttt{P\_L}, \texttt{P\_R}; N02 \texttt{PL44}, \texttt{PR44} \\
\(C_{1644}^{-}\), \(C_{1644}^{+}\) & \texttt{Cm1644} \texttt{Cp1644} &
the products \(\gamma^0\gamma^1\gamma^2\gamma^3\) and
\(\gamma^4\gamma^5\gamma^6\gamma^7\) of the 4+4 matrices;
\(C^{\mp}\gamma^A(C^{\mp})^{-1}=\mp(\gamma^A)^T\) & N01, N02, N10--N12
\texttt{σ16} (\(C^{-}\)); N02 \texttt{S44} (\(C^{+}\)) \\
\(W_{1644}\) & \texttt{W1644} & \(i\,C_{1644}^{-}\gamma_{1644}^4\), the
grading: Hermitian, \(W^2=1\); density of the covariant current & new \\
\(K_{1644}\) & \texttt{K1644} & generator of the evolution of dirac1644
in \(x4\): \(\partial_4\psi=K\psi\) & new \\
\(\beta_{1644}\), \(\alpha_{1644}^j\) & \texttt{beta1644}
\texttt{alpha1644} & the Hermitian matrices \(-i\gamma_{1644}^4\) and
\(-\gamma_{1644}^4\gamma_{1644}^j\) (\(j=0,1,2,3\)) of the quantized 4+4
theory & new \\
\(\mathsf h_{1644}\) & \texttt{h1644} & Hamiltonian matrix of the
quantized 4+4 theory, \(iK_{1644}\) for fields independent of
\(x5,x6,x7\) & new \\
\(\Pi_{1644}\) & \texttt{Pi1644} & the projector
\((i\gamma_{1644}^Ak_A+m_{1644})/2m_{1644}\) on the plane-wave solutions
of the worked example & new \\
\end{longtable}

\textbf{3 Matrices of the 7+1 spacetime (tag 1671)}

\begin{longtable}[]{@{}
  >{\raggedright\arraybackslash}p{(\columnwidth - 6\tabcolsep) * \real{0.1605}}
  >{\raggedright\arraybackslash}p{(\columnwidth - 6\tabcolsep) * \real{0.1235}}
  >{\raggedright\arraybackslash}p{(\columnwidth - 6\tabcolsep) * \real{0.5309}}
  >{\raggedright\arraybackslash}p{(\columnwidth - 6\tabcolsep) * \real{0.1852}}@{}}
\toprule\noalign{}
\begin{minipage}[b]{\linewidth}\raggedright
symbol
\end{minipage} & \begin{minipage}[b]{\linewidth}\raggedright
names in programs and data files
\end{minipage} & \begin{minipage}[b]{\linewidth}\raggedright
meaning
\end{minipage} & \begin{minipage}[b]{\linewidth}\raggedright
origin
\end{minipage} \\
\midrule\noalign{}
\endhead
\bottomrule\noalign{}
\endlastfoot
\(\eta_{1671}\) & \texttt{eta1671} & the matrix of the squares of the
Dirac matrices of 7+1,
\(\gamma_{1671}^A\gamma_{1671}^B+\gamma_{1671}^B\gamma_{1671}^A=2\,\eta_{1671}^{AB}\):
\(\mathrm{diag}(-1,-1,-1,-1,+1,-1,-1,-1)\), time at index 4. The
flat-space part of the project uses it as the metric (time-like
\(=+1\)). It is minus the flat metric \(\eta^{(7+1)}\) of the author's
sign convention for the deflating spacetimes (space-like \(=+1\)) & N02
\texttt{η17} with the index values 0 and 4 exchanged; NOT
\texttt{η71} \\
\(\tau_{1671}^A\) & \texttt{tau1671} & the real \(8\times8\) matrices of
the 7+1 construction: \(\tau_{1671}^4=1_8\), the other seven represent
the imaginary octonion units & N02 \texttt{τO{[}swap04(A){]}} \\
\(\gamma_{1671}^A\) & \texttt{gamma1671} & the eight real \(16\times16\)
Dirac matrices of the 7+1 spacetime (contravariant index) & N02
\texttt{Γ17up{[}swap04(A){]}}; NOT \texttt{Γ71{[}A{]}} \\
\(\Sigma_{1671}^{AB}\) & \texttt{Sigma1671} &
\(\tfrac14[\gamma_{1671}^A,\gamma_{1671}^B]\), the scaled commutators;
they span \(\mathfrak{so}(7,1)\) & N02 \texttt{Σ17{[}A,B{]}}, built from
the \texttt{Γ17} with lower index:
\(\Sigma_{1671}^{AB}=\eta_{1671}^{AA}\eta_{1671}^{BB}\,\)\texttt{Σ17{[}A\textquotesingle{},B\textquotesingle{}{]}}
with \(A'=\mathrm{swap}_{04}(A)\), \(B'=\mathrm{swap}_{04}(B)\) \\
\(\omega_{1671}\) & \texttt{omega1671} & the product
\(\gamma_{1671}^0\cdots\gamma_{1671}^7\); real, \(\omega_{1671}^2=-1\) &
N02 \texttt{Γ17{[}8{]}}, \texttt{ω17} \\
\(\chi_{1671}\) & \texttt{chi1671} & \(i\,\omega_{1671}\), the (complex)
chirality matrix of 7+1 & N02 \texttt{χ17}; NOT \texttt{χ71} \\
\(P_{1671}^{L}\), \(P_{1671}^{R}\) & \texttt{PL1671} \texttt{PR1671} &
the complex chiral projectors \(\tfrac12(1\mp\chi_{1671})\) & N02
\texttt{PL17}, \texttt{PR17}; NOT \texttt{PL71}, \texttt{PR71} \\
\(U_{1671}^{c}\) & \texttt{Uc1671} & the unitary matrix
\((1-i\gamma_{1671}^4)/\sqrt2\) that carries the complex projectors of
7+1 into the real projectors of 4+4 & new \\
\(C_{1671}^{+}\), \(C_{1671}^{-}\) & \texttt{Cp1671} \texttt{Cm1671} &
\(\gamma_{1671}^4\) and \(\gamma_{1671}^4\omega_{1671}\);
\(C^{\pm}\gamma^A(C^{\pm})^{-1}=\pm(\gamma^A)^T\) & new \\
\(\beta_{1671}\), \(\alpha_{1671}^j\) & \texttt{beta1671}
\texttt{alpha1671} & the real symmetric matrices \(\gamma_{1671}^4\) and
\(\gamma_{1671}^4\gamma_{1671}^j\) (seven space indices \(j\)) of the
Hamiltonian & N02 \texttt{βD}, \texttt{αD{[}swap04(j){]}} \\
\(\mathsf h_{1671}\) & \texttt{h1671} & Hamiltonian matrix of the 7+1
theory, \(\alpha_{1671}^jp^j+\beta_{1671}m_{1671}\) & N02 \texttt{h} \\
\(\Pi_{1671}\) & \texttt{Pi1671} & the projector
\((\gamma_{1671}^Ak_A+m_{1671})/2m_{1671}\) of the worked example &
new \\
\(E_{1671}^{L}\), \(E_{1671}^{R}\) & \texttt{EL1671} \texttt{ER1671} &
the chiral frames: \(16\times8\) matrices whose columns span the two
eigenspaces of \(\chi_{1671}\) & N02 \texttt{EL}, \texttt{ER} \\
\(M_{1671}\), \(U_{1671}^{+}\), \(V_{1671}^{-}\) & \texttt{M1671}
\texttt{Upos1671} \texttt{Vneg1671} & the \(8\times8\) matrix
\(m_{1671}+p^j\tau_{1671}^j\) and the matrices of positive- and
negative-energy spinors of the worked example (check \texttt{71-D8}) &
N02 \texttt{M} of its Sec. 7.4 (in N02 the letter \texttt{M} also
denotes other matrices and a mass); the two spinor matrices are new
names \\
\end{longtable}

\textbf{4 Fields and their equations}

\begin{longtable}[]{@{}
  >{\raggedright\arraybackslash}p{(\columnwidth - 6\tabcolsep) * \real{0.1605}}
  >{\raggedright\arraybackslash}p{(\columnwidth - 6\tabcolsep) * \real{0.1235}}
  >{\raggedright\arraybackslash}p{(\columnwidth - 6\tabcolsep) * \real{0.5309}}
  >{\raggedright\arraybackslash}p{(\columnwidth - 6\tabcolsep) * \real{0.1852}}@{}}
\toprule\noalign{}
\begin{minipage}[b]{\linewidth}\raggedright
symbol
\end{minipage} & \begin{minipage}[b]{\linewidth}\raggedright
names in programs and data files
\end{minipage} & \begin{minipage}[b]{\linewidth}\raggedright
meaning
\end{minipage} & \begin{minipage}[b]{\linewidth}\raggedright
origin
\end{minipage} \\
\midrule\noalign{}
\endhead
\bottomrule\noalign{}
\endlastfoot
\(\psi_{1644}\) & \texttt{dirac1644} & the field dirac1644: 16 complex
Grassmann components on the 4+4 spacetime;
\((\gamma_{1644}^A\partial_A-m_{1644})\psi_{1644}=0\) & new;
\texttt{Ψ16} of N01 and N10--N12 is a column of 16 real quantities (in
N10--N12 functions of \texttt{x0}, \texttt{x4} called the wave function
of the universe), not this field \\
\(\psi_{1644}^{L}\), \(\psi_{1644}^{R}\) & \texttt{psiL1644}
\texttt{psiR1644} & the two 8-component chiral parts of dirac1644 (upper
and lower eight components) & new \\
\(m_{1644}\) & \texttt{m1644} & mass of the field dirac1644 & new \\
\(\tilde\psi\) & --- & the covariant adjoint
\(\psi^\dagger C_{1644}^{-}\) of the 4+4 theory (invariant under
\(\mathrm{Spin}_0(4,4)\)) & new \\
\(j_{1644}^A\) & \texttt{j1644} & the conserved current
\(i\,\tilde\psi\gamma_{1644}^A\psi\) of the covariant 4+4 theory
(fermion number) & new \\
\(\psi_{1671}\) & \texttt{dirac1671} & the field dirac1671: 16 complex
Grassmann components on the 7+1 spacetime;
\((i\gamma_{1671}^A\partial_A-m_{1671})\psi_{1671}=0\) & N02
\texttt{dirac16complex} \\
\(\phi_{1671}^{L}\), \(\phi_{1671}^{R}\) & \texttt{phiL1671}
\texttt{phiR1671} & the two 8-component chiral parts of dirac1671 & N02
\texttt{φL}, \texttt{φR} \\
\(m_{1671}\) & \texttt{m1671} & mass of the field dirac1671 & N02
\texttt{m}, \texttt{mD} \\
\(\bar\psi\) & --- & the adjoint \(\psi^\dagger\beta\) of the quantized
theories (\(\beta=\beta_{1671}\) or \(\beta_{1644}\)); in 4+4 it differs
from \(\tilde\psi\) & new \\
\(\beta\), \(\alpha^j\), \(\chi\), \(\mathsf h\), \(m\) & --- & in the
sections common to both spacetimes: the matrices and the mass of the
spacetime at hand, i.e.~the tagged symbols with the tag dropped;
\(\chi\equiv\chi_{1671}\) in 7+1 and \(\chi\equiv\omega_{1644}\) in 4+4
& new \\
\(U(\vec p)\) & --- & the \(16\times8\) matrix of orthonormal
positive-energy spinors at momentum \(\vec p\) & new \\
\end{longtable}

\textbf{5 Symmetries, pairs and order parameter}

\begin{longtable}[]{@{}
  >{\raggedright\arraybackslash}p{(\columnwidth - 6\tabcolsep) * \real{0.1605}}
  >{\raggedright\arraybackslash}p{(\columnwidth - 6\tabcolsep) * \real{0.1235}}
  >{\raggedright\arraybackslash}p{(\columnwidth - 6\tabcolsep) * \real{0.5309}}
  >{\raggedright\arraybackslash}p{(\columnwidth - 6\tabcolsep) * \real{0.1852}}@{}}
\toprule\noalign{}
\begin{minipage}[b]{\linewidth}\raggedright
symbol
\end{minipage} & \begin{minipage}[b]{\linewidth}\raggedright
names in programs and data files
\end{minipage} & \begin{minipage}[b]{\linewidth}\raggedright
meaning
\end{minipage} & \begin{minipage}[b]{\linewidth}\raggedright
origin
\end{minipage} \\
\midrule\noalign{}
\endhead
\bottomrule\noalign{}
\endlastfoot
\(U(1)_F\) & --- & the phase symmetry \(\psi\to e^{i\epsilon}\psi\):
fermion number & new \\
\(U(1)_{VS}\) & \texttt{UVS} & the chiral symmetry
\(\psi\to e^{i\epsilon\chi}\psi\),
\(\Phi_{VS}\to e^{-2i\epsilon}\Phi_{VS}\), broken by the vacuum value of
the Vorton-Scalar & new \\
\(SU(2)_A\), \(SU(2)_B\) & --- & the two factors of the internal
symmetry of the fields that do not depend on the four extra coordinates
(zero-mode sector), in both spacetimes;
\(SO(4)=[SU(2)_A\times SU(2)_B]/Z_2\) & new \\
\(SU(2)_X\), \(SU(2)_Y\) & --- & the two factors of the rotation group
of the four space coordinates \(x0,\dots,x3\) of the 4+4 spacetime (NOT
the internal \(SU(2)_A\times SU(2)_B\)) & new \\
\(SO(3)_T\), \(SU(2)_T\) & --- & the rotations of the three additional
times \(x5,x6,x7\) of 4+4, and its covering group; the diagonal subgroup
of \(SU(2)_A\times SU(2)_B\) & new \\
\(U(1)_E\), \(P\) (generator) & --- & the second central phase symmetry
of the zero-mode sector and its traceless generator, \(P^2=1\): \(P=+1\)
on the doublets of \(SU(2)_A\), \(-1\) on those of \(SU(2)_B\) & new \\
\(\Delta_{1644}^{A}\), \(\Delta_{1671}^{A}\) & \texttt{DeltaV1644}
\texttt{DeltaV1671} & the vector pair operators
\(\psi^TC^{-}\gamma^A\psi\) of each spacetime & new \\
\(\Delta_{1671}^{AB}\) & \texttt{DeltaT1671} & the two-form pair
operators \(\psi^TC^{+}\gamma^A\gamma^B\psi\) of the 7+1 spacetime &
new \\
\(\Delta^a\) & --- & order parameter: expectation value of the vector
pair operator along the extra axis \(a\in\{0,5,6,7\}\);
\(\lvert\Delta\rvert^2\equiv\sum_a\lvert\Delta^a\rvert^2\) & new \\
\(\Phi_X(\vec p)\) & --- & two-particle content
\(U(\vec p)^TXU(-\vec p)\) of the pair operator \(\psi^TX\psi\) & new \\
\(\mathsf M\) & \texttt{M\_gap} & the \(8\times8\) gap matrix on the
internal states of a pair (sans-serif; \texttt{M\_gap} in the statements
of the programs) & new \\
\(d_\pm\) & --- & the two fourfold gaps of a general state,
\(d_\pm^2=\lvert\Delta\rvert^2\pm\sqrt{\lvert\Delta\rvert^4-\lvert\eta_{ab}\Delta^a\Delta^b\rvert^2}\)
& new \\
\(n^a\) & --- & real unit vector: direction of the order parameter in
the space of the four extra axes & new \\
\end{longtable}

\textbf{6 The Vorton-Scalar (suffix VS)}

\begin{longtable}[]{@{}
  >{\raggedright\arraybackslash}p{(\columnwidth - 6\tabcolsep) * \real{0.1605}}
  >{\raggedright\arraybackslash}p{(\columnwidth - 6\tabcolsep) * \real{0.1235}}
  >{\raggedright\arraybackslash}p{(\columnwidth - 6\tabcolsep) * \real{0.5309}}
  >{\raggedright\arraybackslash}p{(\columnwidth - 6\tabcolsep) * \real{0.1852}}@{}}
\toprule\noalign{}
\begin{minipage}[b]{\linewidth}\raggedright
symbol
\end{minipage} & \begin{minipage}[b]{\linewidth}\raggedright
names in programs and data files
\end{minipage} & \begin{minipage}[b]{\linewidth}\raggedright
meaning
\end{minipage} & \begin{minipage}[b]{\linewidth}\raggedright
origin
\end{minipage} \\
\midrule\noalign{}
\endhead
\bottomrule\noalign{}
\endlastfoot
\(\Phi_{VS}\) & \texttt{PhiVS} & the Vorton-Scalar: a complex scalar
field & new (the name is that of the author's project NCSS, where the
scalar is real) \\
\(v_{VS}\) & \texttt{vVS} & vacuum value:
\(\langle\Phi_{VS}\rangle=v_{VS}/\sqrt2\) & new \\
\(\lambda_{VS}\) & \texttt{lambdaVS} \texttt{lamVS} & self-coupling of
the Vorton-Scalar & new \\
\(y_{VS}\) & \texttt{yVS} & Yukawa coupling of the Vorton-Scalar to the
fermion field & new \\
\(s_{VS}\) & \texttt{sVS} & the radial mode of the Vorton-Scalar (real
field) & new \\
\(m_s\) & \texttt{msVS} & mass of the radial mode,
\(\sqrt{2\lambda_{VS}}\,v_{VS}\) & new \\
\(\vartheta_{VS}\) & \texttt{thetaVS} \texttt{thVS} & the phase mode of
the Vorton-Scalar (real field, massless) & new \\
\(g_{VS}\) & \texttt{gVS} & coupling of both modes of the Vorton-Scalar
to the fermions, \(g_{VS}=y_{VS}/\sqrt2=m_Y/v_{VS}\), where \(m_Y\) is
the part of the fermion mass that the Vorton-Scalar gives (equal to
\(m/v_{VS}\) when there is no bare Dirac mass) & new \\
\(\alpha_{VS}\) & \texttt{alphaVS} \texttt{aVS} & \(g_{VS}^2/4\pi\) &
new \\
\(D_{VS}\) & \texttt{DVS} & \(\alpha_{VS}\,m/m_s\): two fermions bind if
and only if \(D_{VS}>D_c\) & new \\
\(J_{VS}^A\) & \texttt{JVS} & the neutral current of the Vorton-Scalar
(bilinear in \(\Phi_{VS}\)) & new \\
\(J_5^A\) & \texttt{J5} & the axial current \(\bar\psi\gamma^A\chi\psi\)
of the fermion field & new \\
\(G_{VS}\) & \texttt{GVS} & constant of the Fermi type of the
neutral-current interaction, \(G_{VS}/\sqrt2=1/(2v_{VS}^2)\) & new \\
\(\mathcal L^{VS}_{1644}\), \(\mathcal L^{VS}_{1671}\) & \texttt{LVS} &
the Lagrangians that define the Vorton-Scalar in the two spacetimes &
new \\
\(C\) (coupling) & --- & \(g_{VS}^2m^2/m_s^2=g_{VS}^4/2\lambda_{VS}\),
the one number on which the fermion mass in matter depends & new \\
\end{longtable}

\textbf{7 Pairing, gap equation and CJT effective action}

\begin{longtable}[]{@{}
  >{\raggedright\arraybackslash}p{(\columnwidth - 6\tabcolsep) * \real{0.1605}}
  >{\raggedright\arraybackslash}p{(\columnwidth - 6\tabcolsep) * \real{0.1235}}
  >{\raggedright\arraybackslash}p{(\columnwidth - 6\tabcolsep) * \real{0.5309}}
  >{\raggedright\arraybackslash}p{(\columnwidth - 6\tabcolsep) * \real{0.1852}}@{}}
\toprule\noalign{}
\begin{minipage}[b]{\linewidth}\raggedright
symbol
\end{minipage} & \begin{minipage}[b]{\linewidth}\raggedright
names in programs and data files
\end{minipage} & \begin{minipage}[b]{\linewidth}\raggedright
meaning
\end{minipage} & \begin{minipage}[b]{\linewidth}\raggedright
origin
\end{minipage} \\
\midrule\noalign{}
\endhead
\bottomrule\noalign{}
\endlastfoot
\(\nu\) & --- & the number of internal states of a particle of given
momentum, \(\nu=8\) & new \\
\(F_S\), \(F_P\) & \texttt{F\_scalar} \texttt{F\_pseudoscalar} &
exchange amplitudes of the radial and of the phase mode in a pair
channel & new \\
\(V_{\rm NC}\), \(V_s(r)\) & --- & contact interaction of the neutral
current in the hidden-axis channel; Yukawa potential of the radial mode
& new \\
\(\sigma_X\) & --- & the sign \(+1\) (vector pairs along the extra axes)
or \(-1\) (their two-form pairs) of the phase-mode term of the kernel &
new \\
\(W(k,k')\) & --- & \(s\)-wave average of the exchange attraction:
kernel of the gap equation & new \\
\(\Delta(k)\), \(\Delta_F\) & --- & gap function; its value at the Fermi
momentum & new \\
\(\xi_k\), \(E_k\), \(\kappa(k)\) & --- & \(\sqrt{m^2+k^2}-\mu\);
quasi-particle energy \(\sqrt{\xi_k^2+\Delta(k)^2}\); pair correlation
\(\langle b(-\vec k)b(\vec k)\rangle\) & new \\
\(k_F\), \(E_F\), \(\mu\) & \texttt{kF} \texttt{kF\_total} \texttt{mu} &
Fermi momentum; Fermi kinetic energy \(\sqrt{m^2+k_F^2}-m\); chemical
potential. The number density is \(n=\nu V_Dk_F^D/(2\pi)^D\) in \(D\)
space dimensions (\(n=\nu k_F^3/6\pi^2\) for \(D=3\)) & new \\
\(N(0)\) & --- & density of states at the Fermi surface per internal
state, \(k_F\sqrt{m^2+k_F^2}/2\pi^2\) & new \\
\(\lambda\), \(\lambda_S\), \(\lambda_P\) & --- & dimensionless pairing
coupling \(N(0)W(k_F,k_F)\) and its parts from the radial and the phase
mode & new \\
\(T_c\) & --- & critical temperature of the paired state & new \\
\(D_c\) & --- & critical value of \(D_{VS}\) for a bound pair in the
Yukawa potential & new \\
\(a_{\rm sc}\), \(a_{BB}\) & --- & scattering length of two fermions;
scattering length of two bound pairs & new \\
\(\Gamma[\phi,G]\), \(\Gamma_2\) & --- & the effective action for
composite operators of Cornwall, Jackiw and Tomboulis; its
two-particle-irreducible part & new (method of the cited paper) \\
\(\mathcal S\), \(D_s\), \(D_\vartheta\) & --- & propagator of the
doubled fermion field; propagators of the radial and of the phase mode &
new \\
\(V[\Delta]\) & --- & the CJT effective potential of the paired state,
measured from the unpaired state & new \\
\(\eta(T)\) & --- & largest eigenvalue of the linearized gap equation,
\(\eta(T_c)=1\) & new \\
\(A_2(T)\), \(A_4\) & --- & coefficients of the expansion of the CJT
effective potential \(V\) along a gap function of fixed shape & new \\
\(a(T)\), \(b\), \(K_{\rm GL}\) & --- & Ginzburg-Landau coefficients of
the weak-coupling free energy (quadratic, quartic and gradient term) &
new \\
\(c_s\) & --- & speed of first sound of the superfluid,
\(c_s^2=dP/d\rho\) & new \\
\end{longtable}

\textbf{8 Matter, temperature and vacuum}

\begin{longtable}[]{@{}
  >{\raggedright\arraybackslash}p{(\columnwidth - 6\tabcolsep) * \real{0.1605}}
  >{\raggedright\arraybackslash}p{(\columnwidth - 6\tabcolsep) * \real{0.1235}}
  >{\raggedright\arraybackslash}p{(\columnwidth - 6\tabcolsep) * \real{0.5309}}
  >{\raggedright\arraybackslash}p{(\columnwidth - 6\tabcolsep) * \real{0.1852}}@{}}
\toprule\noalign{}
\begin{minipage}[b]{\linewidth}\raggedright
symbol
\end{minipage} & \begin{minipage}[b]{\linewidth}\raggedright
names in programs and data files
\end{minipage} & \begin{minipage}[b]{\linewidth}\raggedright
meaning
\end{minipage} & \begin{minipage}[b]{\linewidth}\raggedright
origin
\end{minipage} \\
\midrule\noalign{}
\endhead
\bottomrule\noalign{}
\endlastfoot
\(\phi\) & --- & expectation value of the Vorton-Scalar,
\(\langle\Phi_{VS}\rangle=\phi/\sqrt2\); the fermion mass is
\(g_{VS}\phi\) & new \\
\(m^*\), \(x\) & \texttt{M} \texttt{M\_ks} \texttt{M\_cjt}
\texttt{M\_cjt\_global} \texttt{M\_plus} \texttt{M\_minus}
\texttt{M\_relaxed} \texttt{x} & \(m^*\): the mass of the fermions in
matter; \(x=\phi/v_{VS}\): the expectation value of the Vorton-Scalar in
units of its vacuum value. With a contact interaction \(m^*=m+G_S\,n_S\)
(gap equation of the Kohn-Sham and CJT calculations); with the
Vorton-Scalar and a bare Dirac mass \(m^*=m_D+m_Y\,x+G_S\,n_S\), so that
\(x=m^*/m\) if the whole mass comes from the Vorton-Scalar and no
contact coupling acts. \(m^*_\pm\): the masses of the particles of the
two halves \(W_{1644}=\pm1\) of the 4+4 field in the covariant form of
the self-interaction. The data files and the statements of the programs
write \(m^*\) as \texttt{M}; the author's document KOHN\_SHAM\_GRAVITY
writes \(M\) & new \\
\(\varepsilon(k_F,x)\), \(V_{\rm sea}(x)\) & --- & energy density of
matter as a function of the Fermi momentum and of \(x=m^*/m\) (in units
of \(m^4\) in ordinary space); renormalized energy of the filled
negative-energy states. \(\varepsilon\) without arguments is the energy
density of matter in the units of the calculation at hand & new \\
\(B\) & --- & energy density of the state without expectation value of
the Vorton-Scalar (bag constant) & new \\
\(E_x\) & --- & exchange (Fock) energy of the two modes & new \\
\(M_s\), \(M_\vartheta\), \(a_s\), \(a_\vartheta\) & --- & thermal
masses of the two boson modes in the Hartree approximation, and the
curvatures \(U_{\rm eff}''\), \(U_{\rm eff}'/\phi\) & new \\
\(P(M)\), \(Q(M)\) & --- & thermal part of the boson loop and of
\(\tfrac12\mathrm{Tr}\ln\) & new \\
\(T_{\rm high}\), \(T_N\), \(T_{1/2}\), \(T_0\) & --- & two closed
estimates of the temperature at which the fermion mass vanishes;
temperatures at which the expectation value has fallen to half and to
zero & new \\
\(\sigma(p)\) & --- & self-energy that a pair condensate in the vacuum
gives to the fermions: mass splitting \(m\pm\sigma\) & new \\
\(\alpha_c\), \(\Lambda\) (cutoff) & \texttt{Lambda} & critical value of
\(\alpha_{VS}\) for a pair condensate in the vacuum; \(\Lambda\): a
momentum cutoff (of the ladder equation in the vacuum, and of the filled
states of negative energy in the calculations with contact interaction).
\(\Lambda_8\) and \(\bar\Lambda\) are cosmological terms and have their
own entries & new \\
\end{longtable}

\textbf{9 Cosmology}

\begin{longtable}[]{@{}
  >{\raggedright\arraybackslash}p{(\columnwidth - 6\tabcolsep) * \real{0.1605}}
  >{\raggedright\arraybackslash}p{(\columnwidth - 6\tabcolsep) * \real{0.1235}}
  >{\raggedright\arraybackslash}p{(\columnwidth - 6\tabcolsep) * \real{0.5309}}
  >{\raggedright\arraybackslash}p{(\columnwidth - 6\tabcolsep) * \real{0.1852}}@{}}
\toprule\noalign{}
\begin{minipage}[b]{\linewidth}\raggedright
symbol
\end{minipage} & \begin{minipage}[b]{\linewidth}\raggedright
names in programs and data files
\end{minipage} & \begin{minipage}[b]{\linewidth}\raggedright
meaning
\end{minipage} & \begin{minipage}[b]{\linewidth}\raggedright
origin
\end{minipage} \\
\midrule\noalign{}
\endhead
\bottomrule\noalign{}
\endlastfoot
\(N\) & --- & density of the dark radiation today in units of the
density of one massless neutrino species & new \\
\(z_t\), \(T_d\) & --- & transition redshift of the stepped dark fluid,
\(m=T_{d0}(1+z_t)\); temperature of the dark sector & new \\
\(r_{\rm step}\) & --- & factor by which the density of the dark
radiation rises when the fermions annihilate, \(15^{1/3}\) for this
model & new \\
\(\omega_b\), \(\omega_c\), \(h\), \(H_0\), \(A_s\), \(n_s\), \(\tau\) &
\texttt{tau\_reio} & the six parameters of the standard cosmological
model: densities of baryons and of cold dark matter (\(\Omega h^2\)),
Hubble constant \(H_0=100\,h\) km s\(^{-1}\) Mpc\(^{-1}\), amplitude and
spectral index of the primordial fluctuations, optical depth of
reionization (\texttt{tau\_reio} in the data files; the letter \(\tau\)
with the index \(d\) is the deflation clock) & standard notation \\
\(\Delta\chi^2\) & --- & worsening of the best fit to the mock data
relative to the standard model & new \\
\(\theta_s\), \(r_s\), \(D_A\) & --- & angular size of the sound
horizon; sound horizon; distance to last scattering & standard
notation \\
\(c_{s0}\) & --- & sound speed of superfluid dark matter today at the
mean cosmic density & new \\
\(\rho_0\), \(R_{TF}\) & --- & constant of the equation of state of a
condensate of bound pairs; Thomas-Fermi radius of a core supported by
its pressure & new \\
\(m_{\rm deg}\), \(m_{B,\max}\) & --- & largest fermion mass for
degeneracy, \(m_{\rm deg}=[6\pi^2\hbar^3\rho/(\nu\sigma_v^3)]^{1/4}\),
and largest mass \(m_B\) of a boson for condensation,
\(m_{B,\max}=[(2\pi)^{3/2}\hbar^3\rho/(2.612\,\sigma_v^3)]^{1/4}\), in a
halo of the mass density \(\rho\) and the velocity dispersion
\(\sigma_v\) & new \\
\(\sigma_v\); \(\kappa_{\rm SI}\) & --- & the velocity dispersion of the
dark matter of a halo (with the index \(v\); \(\sigma_8\) is the
amplitude of clustering and \(\sigma(p)\) a self-energy); the bound on
the cross-section per mass of dark matter from colliding clusters, about
\(1~\mathrm{cm^2/g}\) (with the index SI; \(\kappa_8\) is the
gravitational coupling) & new \\
\end{longtable}

\textbf{D1 Deflating spacetimes: hypotheses and geometry}

\begin{longtable}[]{@{}
  >{\raggedright\arraybackslash}p{(\columnwidth - 6\tabcolsep) * \real{0.1605}}
  >{\raggedright\arraybackslash}p{(\columnwidth - 6\tabcolsep) * \real{0.1235}}
  >{\raggedright\arraybackslash}p{(\columnwidth - 6\tabcolsep) * \real{0.5309}}
  >{\raggedright\arraybackslash}p{(\columnwidth - 6\tabcolsep) * \real{0.1852}}@{}}
\toprule\noalign{}
\begin{minipage}[b]{\linewidth}\raggedright
symbol
\end{minipage} & \begin{minipage}[b]{\linewidth}\raggedright
names in programs and data files
\end{minipage} & \begin{minipage}[b]{\linewidth}\raggedright
meaning
\end{minipage} & \begin{minipage}[b]{\linewidth}\raggedright
origin
\end{minipage} \\
\midrule\noalign{}
\endhead
\bottomrule\noalign{}
\endlastfoot
\(\mathrm{g55}\), \(\mathrm{g66}\), \(\mathrm{g77}\) & \texttt{g55}
\texttt{g66} \texttt{g77} & the three rates of deflation of \(x5\),
\(x6\), \(x7\) (positive constants, dimension 1/length). These are the
author's names; they are written upright and never with a subscript.
They are \emph{not} components of the metric: the components are
\(g_{55}\), \(g_{66}\), \(g_{77}\) & author \\
\(\mathrm{rate}_t\) & \texttt{rate\_t} \texttt{rate} & the rate of the
coordinate \(xt\) in a statement that holds for \(t=5,6,7\) alike:
\(\mathrm{rate}_5=\mathrm{g55}\), \(\mathrm{rate}_6=\mathrm{g66}\),
\(\mathrm{rate}_7=\mathrm{g77}\) & new \\
\(\mathrm{gsum}\) & \texttt{gsum} \texttt{gam} &
\(\mathrm{g55}+\mathrm{g66}+\mathrm{g77}\), the total rate: the proper
volume of the three deflating coordinates falls as
\(e^{-\mathrm{gsum}\,x4}\) (\texttt{gam}: gsum in units of 100 km/s/Mpc)
& new \\
\(\mathrm{gsq}\), \(\mathrm{gprod}\) & \texttt{gsq} \texttt{gprod} &
\(\mathrm{g55}^2+\mathrm{g66}^2+\mathrm{g77}^2\) and
\(\mathrm{g55}\,\mathrm{g66}+\mathrm{g55}\,\mathrm{g77}+\mathrm{g66}\,\mathrm{g77}\);
\(\mathrm{gsum}^2=\mathrm{gsq}+2\,\mathrm{gprod}\) & new \\
\(\mathrm{const5}\), \(\mathrm{const6}\), \(\mathrm{const7}\);
\(\mathrm{const}_t\) & \texttt{const5} \texttt{const6} \texttt{const7}
\texttt{const\_t} \texttt{cst} & the three positive constants of the
author's metric (his word: const): the values of the scale factors at
\(x4=0\) & author \\
\(b_5\), \(b_6\), \(b_7\); \(b_t\) & \texttt{b\_5} \texttt{b\_6}
\texttt{b\_7} \texttt{b\_t} & the scale factors of the deflating
coordinates, \(b_t(x4)=\mathrm{const}_t\,e^{-\mathrm{rate}_t\,x4}\): a
coordinate distance \(dxt\) is the proper length
\(b_t\,\lvert dxt\rvert\). Always written with the index 5, 6, 7 or
\(t\) (the letter \(b\) without index is a Ginzburg-Landau coefficient,
and \(\mathsf b\) in sans-serif type is the common scale factor of the
three deflating coordinates in the Einstein-Lovelock evolution) & new \\
\(V_d\) & \texttt{Vd} & \(b_5b_6b_7=\sqrt{\lvert\det g\rvert}\), the
proper volume of a unit coordinate cell of the deflating coordinates;
\(V_d\propto e^{-\mathrm{gsum}\,x4}\) & new \\
\(\mu,\nu,\rho,\sigma\) (indices) & --- & coordinate indices with the
values \(0,\dots,7\) on the deflating spacetimes (only as sub- or
superscripts; as stand-alone letters \(\mu\) is the chemical potential,
\(\nu\) the number of internal states and \(\lambda\) the pairing
coupling). The capital indices \(A,B,C,D\), also with a prime, label the
orthonormal frame: the frame vector \(A\) points along the coordinate
\(xA\) (as a stand-alone letter \(C\) is the coupling
\(g_{VS}^2m^2/m_s^2\), and \(D\) the number of space dimensions with
momentum). In the statements of the programs the indices are spelled out
inside tensor expressions, as in \texttt{g\^{}(mu\ nu)} & new \\
\(t\), \(s\) (indices) & --- & indices of the three additional
coordinates \(x5,x6,x7\), values 5, 6, 7; \(s\) is a second such index
where two are needed (the matrices \(s_h\) carry the index \(h\)) &
new \\
\(e_A{}^\mu\), \(e^A{}_\mu\) & \texttt{e\_A\^{}mu} & the orthonormal
frame and its inverse:
\(e^A{}_\mu=\mathrm{diag}(1,1,1,1,1,b_5,b_6,b_7)\) & new \\
\(\Gamma^\rho{}_{\mu\nu}\) & \texttt{Gamma} \texttt{christoffel} & the
Christoffel symbols of the deflating metric (three indices;
\(\Gamma[\phi,G]\) with square brackets is the CJT effective action).
The author's document KOHN\_SHAM\_GRAVITY writes the Christoffel symbols
with braces and uses \(\Gamma^A\) for the Dirac matrices, which are
\(\gamma^A\) here & new \\
\(R^\rho{}_{\sigma\mu\nu}\), \(R^\mu{}_\nu\), \(R\) & \texttt{riemann}
\texttt{ricci} \texttt{Ric44} \texttt{Ric55} \texttt{Ric66}
\texttt{Ric77} \texttt{Rscalar} & Riemann tensor, Ricci tensor and
curvature scalar of the deflating metric & new \\
\(\mathcal G^\mu{}_\nu\) & \texttt{Ein00} \texttt{Ein44} \texttt{Ein55}
\texttt{Ein66} \texttt{Ein77} \texttt{einstein\_mixed} & the Einstein
tensor \(R^\mu{}_\nu-\tfrac12R\,\delta^\mu_\nu\) (script letter, so that
it cannot be taken for Newton's constant or for a constant of the Fermi
type) & new \\
\(\mathcal X\); \(\mathcal X_{\rm scale}\), \(\mathcal X_{(s)}\) &
\texttt{killing\_vectors} \texttt{killing\_count} & a Killing vector
field of the deflating metric (script letter; the plain letter \(X\) is
the set of the eight coordinates); the scaling vector
\(\partial_4+\sum_t\mathrm{rate}_t\,xt\,\partial_t\); the three special
Killing vectors of the case of equal rates, \(s=5,6,7\) & new \\
\(\bar x4\); \(\bar x5\), \(\bar x6\), \(\bar x7\) & --- & coordinates
of the case of equal rates in which the line element of \(x4,\dots,x7\)
is a multiple of a flat one:
\(\bar x4=e^{\mathrm{g55}\,x4}/\mathrm{g55}\),
\(\bar xt=\mathrm{const}_t\,xt\) & new \\
\(\mathcal H^\mu{}_\nu\); \(\alpha_{GB}\), \(\Lambda_8\) &
\texttt{Lanczos} \texttt{alpha\_GB} \texttt{Lambda\_8} & the Lanczos
tensor of the Gauss-Bonnet term,
\(\mathcal H^\mu{}_\nu=2\big(R\,R^\mu{}_\nu-2R^\mu{}_\alpha R^\alpha{}_\nu-2R^{\mu\alpha}{}_{\nu\beta}R^\beta{}_\alpha+R^{\mu\alpha}{}_{\beta\gamma}R^{\beta\gamma}{}_{\nu\alpha}\big)-\tfrac12\delta^\mu_\nu\big(R^2-4R^\alpha{}_\beta R^\beta{}_\alpha+R^{\alpha\beta}{}_{\gamma\delta}R^{\gamma\delta}{}_{\alpha\beta}\big)\)
(script letter; \(H\) is the Hubble rate); the coupling of the
Gauss-Bonnet term and the cosmological constant of Einstein-Gauss-Bonnet
gravity in eight dimensions, whose equations without matter are
\(\mathcal G^\mu{}_\nu+\Lambda_8\,\delta^\mu_\nu+\alpha_{GB}\,\mathcal H^\mu{}_\nu=0\)
(metric with positive space-like components). In the author's
Einstein-Lovelock notation \(E^{(1)}=-\mathcal G\),
\(E^{(2)}=\mathcal H\), \(\alpha_2=-\alpha_{GB}\),
\(\bar\Lambda=\Lambda_8\) & new \\
\(K_{AB}\) & \texttt{ksec} & the sectional curvature of the plane of the
frame directions \(A\) and \(B\) (metric with positive space-like
components): on the deflating spacetimes \(K_{4t}=\mathrm{rate}_t^2\),
\(K_{ts}=\mathrm{rate}_t\,\mathrm{rate}_s\) for \(t\neq s\), zero if one
direction is flat; for scale factors that depend on \(x4\) in any way,
\(K_{4i}=H_i'+H_i^2\) and \(K_{ij}=H_iH_j\) with the rate \(H_i\) of the
direction \(i\) (the pair curvatures of the author's Einstein-Lovelock
equations). Always with two subscripts: the letter \(K\) without
subscript is the generator of the evolution of the 4+4 field in \(x4\),
\(K\) as an index is an index set of the self-interaction, and \(K\)
with an order as subscript and with an argument is a modified Bessel
function & new \\
\(\varsigma\) & \texttt{affine} & affine parameter along a geodesic; for
a massive particle on the 7+1 spacetime, the proper time divided by the
mass & new \\
\(P_\mu\) & \texttt{P\_1} \texttt{P\_5} & the conserved covariant
momenta of a geodesic (with an index; \(P\) without index is the
generator of \(U(1)_E\)) & new \\
\(x4_{\rm turn}\) & \texttt{x4\_turn} & the value of the ordinary time
at which a classical particle with momentum along a deflating time turns
back, and beyond which a mode that depends on a deflating time is
extinguished (4+4): there \(k_d=E_\perp\) & new \\
\end{longtable}

\textbf{D2 Deflating spacetimes: matrices, fields and their equations}

\begin{longtable}[]{@{}
  >{\raggedright\arraybackslash}p{(\columnwidth - 6\tabcolsep) * \real{0.1605}}
  >{\raggedright\arraybackslash}p{(\columnwidth - 6\tabcolsep) * \real{0.1235}}
  >{\raggedright\arraybackslash}p{(\columnwidth - 6\tabcolsep) * \real{0.5309}}
  >{\raggedright\arraybackslash}p{(\columnwidth - 6\tabcolsep) * \real{0.1852}}@{}}
\toprule\noalign{}
\begin{minipage}[b]{\linewidth}\raggedright
symbol
\end{minipage} & \begin{minipage}[b]{\linewidth}\raggedright
names in programs and data files
\end{minipage} & \begin{minipage}[b]{\linewidth}\raggedright
meaning
\end{minipage} & \begin{minipage}[b]{\linewidth}\raggedright
origin
\end{minipage} \\
\midrule\noalign{}
\endhead
\bottomrule\noalign{}
\endlastfoot
\(\gamma_{1644d}^\mu\), \(\gamma_{1671d}^\mu\) & \texttt{gamma\_1644d}
\texttt{gamma\_1671d} & the Dirac matrices of the deflating spacetimes
in coordinate components, \(\gamma^\mu=e_A{}^\mu\gamma^A\):
\(\gamma^t=\gamma^{A=t}/b_t\) for \(t=5,6,7\), the other five are the
constant matrices \(\gamma_{1644}^A\), \(\gamma_{1671}^A\) of the frame
& new \\
\(\omega_{\mu AB}\) & \texttt{omega\_(mu} \texttt{A} \texttt{B)} & the
coefficients of the spin connection,
\(\omega_{\mu AB}=\eta_{AC}\,e^C{}_\nu\big(\partial_\mu e_B{}^\nu+\Gamma^\nu{}_{\mu\lambda}e_B{}^\lambda\big)\)
with the flat metric \(\eta=\eta^{(4+4)}\) or \(\eta^{(7+1)}\) (three
indices; \(\omega_{1644}\), \(\omega_{1671}\) with a tag are the
products of all eight Dirac matrices) & new \\
\(\Omega_\mu\) & \texttt{Omega\_mu} \texttt{omega} & the spin connection
as a matrix,
\(\Omega_\mu=\tfrac14\,c_\gamma\,\omega_{\mu AB}\gamma^A\gamma^B\); the
covariant derivative of a spinor is
\(\nabla_\mu\psi=(\partial_\mu+\Omega_\mu)\psi\). On both deflating
spacetimes \(\Omega_t=\tfrac12\,\mathrm{rate}_t\,b_t\,\gamma^t\gamma^4\)
(frame matrices) for \(t=5,6,7\) and \(\Omega_\mu=0\) otherwise.
(\(\Omega_{\rm CJT}\) is an effective potential) & new \\
\(\psi_{1644d}\), \(\psi_{1671d}\) & \texttt{dirac1644d}
\texttt{dirac1671d} & the fields dirac1644d and dirac1671d: the
16-component fields on the deflating spacetimes. With mass term and
self-interaction they obey
\((\gamma^\mu\nabla_\mu-m)\psi_{1644d}=\sum_rg_r\sum_{\lvert K\rvert=r}(\tilde\psi\Theta_K\psi)\Theta^K\psi\)
(covariant form of 4+4) and
\((i\gamma^\mu\nabla_\mu-m)\psi_{1671d}=\sum_rg_r\sum_{\lvert K\rvert=r}(\bar\psi\Theta_K\psi)\Theta^K\psi\)
(7+1) & new \\
\(\hat\psi\) & \texttt{psihat} & the weighted field
\(\sqrt{V_d}\,\psi\), which obeys the equation without spin connection &
new \\
\(K(x4)\), \(\mathsf h(x4)\) & \texttt{K} \texttt{h} & the generator
\(K_{1644}\) and the Hamiltonian matrix \(\mathsf h_{1671}\) with the
wavenumbers along the deflating coordinates divided by the scale
factors: without self-interaction \(\partial_4\hat\psi=K\hat\psi\)
(4+4), \(i\partial_4\hat\psi=\mathsf h\hat\psi\) (7+1) & new \\
\end{longtable}

\textbf{D3 Deflating spacetimes: solutions mode by mode}

\begin{longtable}[]{@{}
  >{\raggedright\arraybackslash}p{(\columnwidth - 6\tabcolsep) * \real{0.1605}}
  >{\raggedright\arraybackslash}p{(\columnwidth - 6\tabcolsep) * \real{0.1235}}
  >{\raggedright\arraybackslash}p{(\columnwidth - 6\tabcolsep) * \real{0.5309}}
  >{\raggedright\arraybackslash}p{(\columnwidth - 6\tabcolsep) * \real{0.1852}}@{}}
\toprule\noalign{}
\begin{minipage}[b]{\linewidth}\raggedright
symbol
\end{minipage} & \begin{minipage}[b]{\linewidth}\raggedright
names in programs and data files
\end{minipage} & \begin{minipage}[b]{\linewidth}\raggedright
meaning
\end{minipage} & \begin{minipage}[b]{\linewidth}\raggedright
origin
\end{minipage} \\
\midrule\noalign{}
\endhead
\bottomrule\noalign{}
\endlastfoot
\(k_5\), \(k_6\), \(k_7\); \(k_t\) & \texttt{k\_5} \texttt{k\_6}
\texttt{k\_7} \texttt{k\_t} & the conserved wavenumbers of a mode along
the deflating coordinates (covariant components of the wave vector
\(k_A\)) & new \\
\(E_\perp\) & \texttt{E\_perp} & the energy of a mode, or of a classical
particle, without its motion along the deflating coordinates:
\(\sqrt{m^2+k_0^2+k_1^2+k_2^2+k_3^2}\) (with the wavenumbers or momenta
along \(x0,\dots,x3\)) & new \\
\(k_d\) & \texttt{k\_d} & the physical wavenumber along the deflating
coordinates, \(k_d(x4)=\big[\sum_tk_t^2/b_t(x4)^2\big]^{1/2}\); it grows
as the coordinates deflate & new \\
\(\mathsf h_\perp\), \(K_\perp\) & \texttt{h\_perp} \texttt{K\_perp} &
the parts of \(\mathsf h\) and of \(K\) that do not contain the scale
factors: the operators of the modes that do not depend on the deflating
coordinates, \(\mathsf h_\perp=m\beta+\sum_ap_a\alpha^a\) with
\(\mathsf h_\perp^2=E_\perp^2\), and
\(K_\perp=-m\gamma^4+i\sum_ak_a\gamma^4\gamma^a\) with
\(K_\perp^2=-E_\perp^2\) & new \\
\(\gamma^d\), \(\alpha^d\) & --- & the frame matrix along the direction
of the wave vector in the deflating coordinates,
\(\gamma^d=k_d^{-1}\sum_t(k_t/b_t)\gamma^t\), and
\(\alpha^d=\gamma^4\gamma^d\); constant in time when one rate governs &
new \\
\(w\); \(f_+\), \(f_-\); \(Z_+\), \(Z_-\) & --- & auxiliary objects of
the exact solutions of the equations without self-interaction: a
constant eigenvector of \(\mathsf h_\perp\) or of \(K_\perp\); the two
coefficient functions of the ansatz; two cylinder functions whose orders
differ by one & new \\
\(u_d\) & \texttt{u\_d} & \(k_d/\mathrm{rate}_t\) for a mode that
depends on one deflating coordinate: the argument of the Bessel
functions of the exact solutions & new \\
\(n_{\rm cr}\) & \texttt{n\_cr} \texttt{n\_min} \texttt{n\_mean}
\texttt{n\_max} \texttt{n\_exact} \texttt{fd5} \texttt{fd6} & the number
of particles per internal state that the deflation creates in a mode
that starts empty (7+1) & new \\
\(T_{\rm defl}\) & --- & the temperature of the thermal spectrum of the
created quanta, \(\mathrm{rate}_t/2\pi\) in units with \(\hbar=k_B=1\)
(7+1) & new \\
\(S_{\rm WKB}\) & \texttt{S\_wkb} & \(\int\sqrt{k_d^2-E_\perp^2}\,dx4\)
from the turning value on: exponent of the extinction of a bounded
solution (4+4) & new \\
\(\varphi_{\rm bos}\), \(\hat\varphi_{\rm bos}\), \(m_{\rm bos}\) &
\texttt{phi\_bos} \texttt{Shat} \texttt{Shat2} \texttt{m\_bos} & a mode
of the Vorton-Scalar (the radial mode, \(m_{\rm bos}=m_s\), or the phase
mode, \(m_{\rm bos}=0\)), its weighted form
\(\sqrt{V_d}\,\varphi_{\rm bos}\), and its mass & new \\
\(\varpi_b\) & \texttt{varpi\_b} &
\(\sqrt{k^2+m_{\rm bos}^2-\mathrm{gsum}^2/4}\): the frequency of the
weighted boson mode that does not depend on the deflating coordinates &
new \\
\end{longtable}

\textbf{D4 The deflation clock and the couplings}

\begin{longtable}[]{@{}
  >{\raggedright\arraybackslash}p{(\columnwidth - 6\tabcolsep) * \real{0.1605}}
  >{\raggedright\arraybackslash}p{(\columnwidth - 6\tabcolsep) * \real{0.1235}}
  >{\raggedright\arraybackslash}p{(\columnwidth - 6\tabcolsep) * \real{0.5309}}
  >{\raggedright\arraybackslash}p{(\columnwidth - 6\tabcolsep) * \real{0.1852}}@{}}
\toprule\noalign{}
\begin{minipage}[b]{\linewidth}\raggedright
symbol
\end{minipage} & \begin{minipage}[b]{\linewidth}\raggedright
names in programs and data files
\end{minipage} & \begin{minipage}[b]{\linewidth}\raggedright
meaning
\end{minipage} & \begin{minipage}[b]{\linewidth}\raggedright
origin
\end{minipage} \\
\midrule\noalign{}
\endhead
\bottomrule\noalign{}
\endlastfoot
\(\tau_d\) & \texttt{tau\_d} & the deflation clock,
\(\tau_d=\mathrm{gsum}\,(x4-x4_{\rm ref})\): the number of e-foldings by
which the volume \(V_d\) has shrunk since the reference time; in an
evolution in which the rates change, \(\tau_d=-\ln(V_d/V_d^{\rm ref})\)
(the letter \(\tau\) without the index \(d\) is the optical depth of
reionization) & new \\
\(\tau_{\rm pair}\) & \texttt{tau\_pair} & the clock measured from the
moment at which two fermions become bound: \(\ln(D_{VS}/D_c)\) & new \\
\(\hat\Phi_{VS}\) & \texttt{Phihat} & the weighted Vorton-Scalar
\(\sqrt{V_d}\,\Phi_{VS}\) & new \\
\end{longtable}

\textbf{D5 Dynamics of the paired state}

\begin{longtable}[]{@{}
  >{\raggedright\arraybackslash}p{(\columnwidth - 6\tabcolsep) * \real{0.1605}}
  >{\raggedright\arraybackslash}p{(\columnwidth - 6\tabcolsep) * \real{0.1235}}
  >{\raggedright\arraybackslash}p{(\columnwidth - 6\tabcolsep) * \real{0.5309}}
  >{\raggedright\arraybackslash}p{(\columnwidth - 6\tabcolsep) * \real{0.1852}}@{}}
\toprule\noalign{}
\begin{minipage}[b]{\linewidth}\raggedright
symbol
\end{minipage} & \begin{minipage}[b]{\linewidth}\raggedright
names in programs and data files
\end{minipage} & \begin{minipage}[b]{\linewidth}\raggedright
meaning
\end{minipage} & \begin{minipage}[b]{\linewidth}\raggedright
origin
\end{minipage} \\
\midrule\noalign{}
\endhead
\bottomrule\noalign{}
\endlastfoot
\(\Xi_d\) & \texttt{Xi\_d} &
\(\mathrm{gsum}/(\lambda\,\Delta_{\rm eq})\): decides whether the paired
state can follow the deflation (\(\Xi_d\ll1\): yes) & new \\
\(n_k\) & --- & occupation of the pair of states of the single-particle
energy \(\xi_k\) in the shell model of the paired state (with the index
\(k\); the letter \(n\) without index is a density) & new \\
\(\Delta_{\rm eq}\) & \texttt{gap\_eq} & the equilibrium gap for the
pairing strength of the moment & new \\
\(E_{\rm sh}\) & \texttt{E\_sh} & half-width in energy of the shell of
states around the Fermi surface in the model of the pair dynamics; unit
of energy and of inverse time there & new \\
\(n_{\rm ex}\) & \texttt{n\_ex} \texttt{n\_ex\_final} & density of
broken pairs (in units of \(E_{\rm sh}\) times the density of states) &
new \\
\end{longtable}

\textbf{D6 Cosmology and bounds on the deflating spacetimes}

\begin{longtable}[]{@{}
  >{\raggedright\arraybackslash}p{(\columnwidth - 6\tabcolsep) * \real{0.1605}}
  >{\raggedright\arraybackslash}p{(\columnwidth - 6\tabcolsep) * \real{0.1235}}
  >{\raggedright\arraybackslash}p{(\columnwidth - 6\tabcolsep) * \real{0.5309}}
  >{\raggedright\arraybackslash}p{(\columnwidth - 6\tabcolsep) * \real{0.1852}}@{}}
\toprule\noalign{}
\begin{minipage}[b]{\linewidth}\raggedright
symbol
\end{minipage} & \begin{minipage}[b]{\linewidth}\raggedright
names in programs and data files
\end{minipage} & \begin{minipage}[b]{\linewidth}\raggedright
meaning
\end{minipage} & \begin{minipage}[b]{\linewidth}\raggedright
origin
\end{minipage} \\
\midrule\noalign{}
\endhead
\bottomrule\noalign{}
\endlastfoot
\(G_N\), \(G_8\) & \texttt{G\_N} \texttt{GN\_early} \texttt{GN\_star}
\texttt{G\_8} & Newton's constant of ordinary space, and the
gravitational coupling of the eight-dimensional Einstein theory
(assumption A2):
\(G_N=G_8/(\text{proper volume of the extra coordinates})\propto e^{\mathrm{gsum}\,x4}\)
& new \\
\(a_{\rm sp}\), \(H\) & --- & scale factor of ordinary space (assumption
A1) and its rate \(H=(da_{\rm sp}/dx4)/a_{\rm sp}\); \(H_0\) is its
value today. (The rates of the Einstein-Lovelock evolution are
\(H_{\mathsf a}\), \(H_{\mathsf b}\); the author's Einstein-Lovelock
notebooks use the letter \(H\) for the rate in the exponent of their
metric, which corresponds here to \(\mathrm{g55}\), \(\mathrm{g66}\),
\(\mathrm{g77}\)) & new \\
\(t_0\) & \texttt{age} \texttt{age\_Gyr} & the age of the universe in
the model at hand & new \\
\(\theta_*\), \(z_*\) & \texttt{theta\_star} \texttt{z\_star} & the
angle under which the sound horizon at last scattering is seen,
\(r_s(z_*)/D_A(z_*)\) with the comoving distance \(D_A\), and the
redshift \(z_*\) of last scattering at which the optical depth without
reionization is 1 (the \(100\,\theta_*\) of the Planck papers). The
angle \(\theta_s\) is the same ratio at the maximum of the visibility
function; for the model without deflation \(100\,\theta_*=1.03965\) and
\(100\,\theta_s=1.040358\) & new \\
\(\delta_m\), \(f_g\), \(\sigma_8\) & \texttt{growth} \texttt{f\_growth}
\texttt{growth\_ratio} \texttt{f\_growth\_ratio}
\texttt{f\_sigma8\_ratio} & density contrast of matter in ordinary
space; its growth rate \(d\ln\delta_m/d\ln a_{\rm sp}\) today; the
amplitude of clustering (only ratios of \(\sigma_8\) are used) & new \\
\(v_{\rm EW}\), \(n_{\rm EW}\) & \texttt{v\_EW} \texttt{n\_EW} & the
electroweak scale (the vacuum value of the Higgs field of the Standard
Model, 246 GeV), and the exponent with which it would follow the
coupling under the hypothesis H2:
\(n_{\rm EW}=d\ln v_{\rm EW}/d\ln\alpha_{VS}\) & new \\
\(\varrho_c\) & \texttt{rho\_c} & abbreviation in the formula for
\(n_{\rm EW}\) of the scale-invariant ladder equation:
\(\varrho_c=\tfrac12\sqrt{\alpha_{VS}/\alpha_c-1}\) (the letter \(\rho\)
without index is a density) & new \\
\(m_p/m_e\) & \texttt{mu\_pe} & the ratio of the proton mass to the
electron mass (always written as this ratio in the documents;
\texttt{mu\_pe} in the statements of the programs) & new \\
\(G_{\rm Fermi}\), \(\alpha_{\rm em}\) & \texttt{G\_Fermi} & the Fermi
constant of the weak interaction and the fine-structure constant (they
enter only in the discussion of bounds) & new \\
\end{longtable}

\textbf{D7 Symmetries and vortex lines on the deflating spacetimes}

\begin{longtable}[]{@{}
  >{\raggedright\arraybackslash}p{(\columnwidth - 6\tabcolsep) * \real{0.1605}}
  >{\raggedright\arraybackslash}p{(\columnwidth - 6\tabcolsep) * \real{0.1235}}
  >{\raggedright\arraybackslash}p{(\columnwidth - 6\tabcolsep) * \real{0.5309}}
  >{\raggedright\arraybackslash}p{(\columnwidth - 6\tabcolsep) * \real{0.1852}}@{}}
\toprule\noalign{}
\begin{minipage}[b]{\linewidth}\raggedright
symbol
\end{minipage} & \begin{minipage}[b]{\linewidth}\raggedright
names in programs and data files
\end{minipage} & \begin{minipage}[b]{\linewidth}\raggedright
meaning
\end{minipage} & \begin{minipage}[b]{\linewidth}\raggedright
origin
\end{minipage} \\
\midrule\noalign{}
\endhead
\bottomrule\noalign{}
\endlastfoot
\(P_c\) & \texttt{Pc} & \(i\alpha^1\alpha^2\alpha^3\chi\), the product
of the five matrices of the sector that does not depend on the extra
coordinates; equal, up to sign, to the generator \(P\) of \(U(1)_E\) &
new \\
\(\theta_{12}\) & --- & the angle around a vortex line that lies along
\(x3\), in the plane of \(x1\) and \(x2\) (with the index 12;
\(\theta_*\) is the angle of the sound horizon and \(\theta(\cdot)\) the
step function) & new \\
\(w_z\) & \texttt{w\_z} & a constant spinor of a fermion zero mode on a
vortex line & new \\
\end{longtable}

\textbf{D8 Tags and identifiers of the deflating spacetimes}

\begin{longtable}[]{@{}
  >{\raggedright\arraybackslash}p{(\columnwidth - 6\tabcolsep) * \real{0.1605}}
  >{\raggedright\arraybackslash}p{(\columnwidth - 6\tabcolsep) * \real{0.1235}}
  >{\raggedright\arraybackslash}p{(\columnwidth - 6\tabcolsep) * \real{0.5309}}
  >{\raggedright\arraybackslash}p{(\columnwidth - 6\tabcolsep) * \real{0.1852}}@{}}
\toprule\noalign{}
\begin{minipage}[b]{\linewidth}\raggedright
symbol
\end{minipage} & \begin{minipage}[b]{\linewidth}\raggedright
names in programs and data files
\end{minipage} & \begin{minipage}[b]{\linewidth}\raggedright
meaning
\end{minipage} & \begin{minipage}[b]{\linewidth}\raggedright
origin
\end{minipage} \\
\midrule\noalign{}
\endhead
\bottomrule\noalign{}
\endlastfoot
1644d, 1671d & \texttt{1644d} \texttt{1671d} & tags of the deflating 4+4
and 7+1 spacetimes and of every object defined on them (the tags 1644,
1671 of the flat spacetimes with the letter d) & new \\
A1, A2, A3, H2, HV & --- & A1, A2: the two assumptions of the
cosmological part (ordinary space expands; gravity is the
eight-dimensional Einstein theory). A3: the assumption about the
perturbations of ordinary space (the scale factors of the extra
coordinates are not perturbed, and the agent that holds them has no
energy density), with its closures. H2: the author's hypothesis that the
superfluid state is the Higgs field. HV: the different assumption that
the vacuum value of the Vorton-Scalar sets the electroweak scale &
new \\
"Momentum", "Pressure", "Shear" (closures) & \texttt{Momentum}
\texttt{Pressure} \texttt{Shear} \texttt{closure} & the three closures
of the perturbation equations under assumption A3, named after the
Einstein equation of ordinary space that holds with matter and radiation
alone besides the energy constraint: the momentum constraint, the trace
(pressure) equation, or the equation of the shear. In data files the
suffixes \texttt{\_M} and \texttt{\_P} mark the first two; \texttt{\_N4}
marks the four-dimensional growth equation, which is none of them &
new \\
\end{longtable}

\textbf{D9 Perturbations of ordinary space and the microwave background}

\begin{longtable}[]{@{}
  >{\raggedright\arraybackslash}p{(\columnwidth - 6\tabcolsep) * \real{0.1605}}
  >{\raggedright\arraybackslash}p{(\columnwidth - 6\tabcolsep) * \real{0.1235}}
  >{\raggedright\arraybackslash}p{(\columnwidth - 6\tabcolsep) * \real{0.5309}}
  >{\raggedright\arraybackslash}p{(\columnwidth - 6\tabcolsep) * \real{0.1852}}@{}}
\toprule\noalign{}
\begin{minipage}[b]{\linewidth}\raggedright
symbol
\end{minipage} & \begin{minipage}[b]{\linewidth}\raggedright
names in programs and data files
\end{minipage} & \begin{minipage}[b]{\linewidth}\raggedright
meaning
\end{minipage} & \begin{minipage}[b]{\linewidth}\raggedright
origin
\end{minipage} \\
\midrule\noalign{}
\endhead
\bottomrule\noalign{}
\endlastfoot
\(\tau_{\rm cf}\) & \texttt{tau\_cf} \texttt{tau0\_Mpc} & the conformal
time of ordinary space, \(d\tau_{\rm cf}=dx4/a_{\rm sp}\), in Mpc
(\texttt{tau0\_Mpc}: its value today). In the equations of the
perturbations a prime is the derivative with respect to
\(\tau_{\rm cf}\) & new \\
\(H_{\rm cf}\), \(H_e\) & \texttt{Hc} \texttt{He} \texttt{hub}
\texttt{hub\_e} \texttt{a\_gsum\_over\_Hc} & the conformal Hubble rate
\(H_{\rm cf}=a_{\rm sp}H\), and
\(H_e=H_{\rm cf}-a_{\rm sp}\,\mathrm{gsum}/2\), which takes its place in
three of the Einstein equations of ordinary space when the extra
coordinates deflate & new \\
\(k\) (perturbations), \(\ell\), \(D_\ell^{TT}\), \(D_\ell^{TE}\),
\(D_\ell^{EE}\) & --- & the comoving wavenumber of a perturbation of
ordinary space (in Mpc\(^{-1}\); the physical wavenumber is
\(k/a_{\rm sp}\)); the multipole; and the spectra of the temperature, of
the polarization and their cross spectrum,
\(D_\ell=\ell(\ell+1)C_\ell/2\pi\) in \(\mu\)K\(^2\) & new \\
\(h_{\rm sy}\), \(\eta_{\rm sy}\), \(\alpha_{\rm sy}\) & --- & the two
functions that describe a scalar perturbation of the metric of ordinary
space in the synchronous gauge (the \(h\) and \(\eta\) of Ma and
Bertschinger; here with the index sy, because \(h\) is the Hubble
parameter and \(\eta\) the flat metric), and the combination
\(\alpha_{\rm sy}=(h_{\rm sy}'+6\eta_{\rm sy}')/2k^2\) & new \\
\(\mathcal U\) & \texttt{potential} &
\(\eta_{\rm sy}+\alpha_{\rm sy}'\), the sum of the two gravitational
potentials of the conformal Newtonian gauge (script letter): the
potential that light feels & new \\
\(\mathcal D\), \(\mathcal T\), \(\mathcal P\), \(\mathcal A\) &
\texttt{d\_sum} \texttt{three\_dpres} & \(8\pi G_N a_{\rm sp}^2\) times
the perturbation of the density, of the divergence of the momentum
density, of the pressure, and of the anisotropic stress of matter and
radiation together (script letters; in the notation of Ma and
Bertschinger: \(\delta\rho\), \((\rho+P)\theta\), \(\delta P\),
\((\rho+P)\sigma\)) & new \\
\(\mathcal E\), \(\mathcal E_P\), \(\mathcal E_m\), \(\mathcal E_r\),
\(\mathcal B\) & \texttt{grho\_sum} \texttt{half\_rho\_plus\_p}
\texttt{b\_trace} \texttt{n\_bg} & \(8\pi G_N a_{\rm sp}^2\) times the
density and the pressure of the unperturbed universe, summed over the
species (\(\mathcal E_m\), \(\mathcal E_r\): the parts of \(\mathcal E\)
from pressureless matter and from radiation); and
\(\mathcal B=H_e'-H_e^2+\tfrac12(\mathcal E+\mathcal E_P)\), which
vanishes without deflation (script letters) & new \\
\(p_{\rm agent}\) & --- & the pressure in ordinary space of the agent,
that is, of whatever keeps the three rates of deflation constant, in the
unperturbed universe (under assumption A2 the agent has no energy
density) & new \\
\(\mathcal Y_M\), \(\mathcal Y_T\), \(\mathcal Y_S\); \(\mathcal Z\) &
\texttt{momentum\_gap} \texttt{trace\_gap} \texttt{shear\_gap}
\texttt{x\_gap} \texttt{shear\_check} \texttt{trace\_check} & the gaps:
what the momentum constraint, the trace equation and the equation of the
shear lack when only matter and radiation stand on their right-hand
sides; the agent that holds the extra scale factors supplies them.
\(\mathcal Z=h_{\rm sy}'\,\mathcal B-\tfrac12a_{\rm sp}\,\mathrm{gsum}\,(\mathcal D-3\mathcal P)\)
(script letters) & new \\
\(Y_{\rm He}\) & --- & the mass fraction of helium after nucleosynthesis
(with the index He; \texttt{Y\_He} in the data files) & new \\
\end{longtable}

\textbf{10 The metric of the deflating spacetimes in the author's sign
convention}

\begin{longtable}[]{@{}
  >{\raggedright\arraybackslash}p{(\columnwidth - 6\tabcolsep) * \real{0.1605}}
  >{\raggedright\arraybackslash}p{(\columnwidth - 6\tabcolsep) * \real{0.1235}}
  >{\raggedright\arraybackslash}p{(\columnwidth - 6\tabcolsep) * \real{0.5309}}
  >{\raggedright\arraybackslash}p{(\columnwidth - 6\tabcolsep) * \real{0.1852}}@{}}
\toprule\noalign{}
\begin{minipage}[b]{\linewidth}\raggedright
symbol
\end{minipage} & \begin{minipage}[b]{\linewidth}\raggedright
names in programs and data files
\end{minipage} & \begin{minipage}[b]{\linewidth}\raggedright
meaning
\end{minipage} & \begin{minipage}[b]{\linewidth}\raggedright
origin
\end{minipage} \\
\midrule\noalign{}
\endhead
\bottomrule\noalign{}
\endlastfoot
\(g^{(4+4)}_{\mu\nu}\), \(g^{(7+1)}_{\mu\nu}\) & \texttt{gs} & the
metrics of the deflating spacetimes in the author's sign convention of
November 13, 2026 (space-like components positive, time-like components
negative):
\(g^{(4+4)}=\mathrm{diag}(+1,+1,+1,+1,-1,-b_5^2,-b_6^2,-b_7^2)\) and
\(g^{(7+1)}=\mathrm{diag}(+1,+1,+1,+1,-1,+b_5^2,+b_6^2,+b_7^2)\), that
is \(g^{(4+4)}_{tt}=-\mathrm{const}_t^2\,e^{-2\,\mathrm{rate}_t\,x4}\)
and \(g^{(7+1)}_{tt}=+\mathrm{const}_t^2\,e^{-2\,\mathrm{rate}_t\,x4}\)
for \(t=5,6,7\). A component of the metric carries the superscript
\((4+4)\) or \((7+1)\) and two coordinate indices; it is not a rate
(\(\mathrm{g55}\)) and not a coupling (\(g_r\), \(g_{VS}\)). The first
revision on the deflating spacetimes wrote \(g_{1644d}=g^{(4+4)}\) and
\(g_{1671d}=-g^{(7+1)}\) (time-like component positive on 7+1); these
two symbols are retired & author \\
\(\eta^{(4+4)}\), \(\eta^{(7+1)}\) & --- & the flat metrics of the
author's sign convention, with which frame indices are raised and
lowered on the deflating spacetimes:
\(\eta^{(4+4)}=\eta_{1644}=\mathrm{diag}(+1,+1,+1,+1,-1,-1,-1,-1)\) and
\(\eta^{(7+1)}=-\eta_{1671}=\mathrm{diag}(+1,+1,+1,+1,-1,+1,+1,+1)\);
\(g_{\mu\nu}\,e_A{}^\mu e_B{}^\nu=\eta_{AB}\) & author \\
\(c_\gamma\) & \texttt{clifford\_sign} & the sign that relates the Dirac
matrices to the metric of the author's sign convention,
\(\gamma^A\gamma^B+\gamma^B\gamma^A=2\,c_\gamma\,\eta^{AB}\) with
\(\eta=\eta^{(4+4)}\) or \(\eta^{(7+1)}\): \(c_\gamma=+1\) on 4+4 and
\(c_\gamma=-1\) on 7+1. The Dirac matrices themselves are those of the
flat-space part of the project, unchanged & new \\
\(b_t^{\rm ref}\) & \texttt{s5} \texttt{s6} \texttt{s7} & the value of
the scale factor \(b_t\) at the reference time \(x4_{\rm ref}\). A
particle that keeps its wavenumber \(k_t\) has the proper momentum
\(k_t/b_t\), larger than at the reference time by the stretch factor
\(b_t^{\rm ref}/b_t=e^{\mathrm{rate}_t(x4-x4_{\rm ref})}\) (columns
\texttt{s5}, \texttt{s6}, \texttt{s7}) & new \\
\(\tau_d^{0}\) & \texttt{tau\_ref} \texttt{tau\_ref\_end} & the reading
that the deflation clock would have if the rates had kept their initial
value, \(\tau_d^0=3\,\mathrm{g55}\,x4\) in the evolutions with gravity
(equal rates), where the clock actually reached is
\(\tau_d=-3\ln\mathsf b\) & new \\
\end{longtable}

\textbf{11 Mass term and self-interaction of the fermion field}

\begin{longtable}[]{@{}
  >{\raggedright\arraybackslash}p{(\columnwidth - 6\tabcolsep) * \real{0.1605}}
  >{\raggedright\arraybackslash}p{(\columnwidth - 6\tabcolsep) * \real{0.1235}}
  >{\raggedright\arraybackslash}p{(\columnwidth - 6\tabcolsep) * \real{0.5309}}
  >{\raggedright\arraybackslash}p{(\columnwidth - 6\tabcolsep) * \real{0.1852}}@{}}
\toprule\noalign{}
\begin{minipage}[b]{\linewidth}\raggedright
symbol
\end{minipage} & \begin{minipage}[b]{\linewidth}\raggedright
names in programs and data files
\end{minipage} & \begin{minipage}[b]{\linewidth}\raggedright
meaning
\end{minipage} & \begin{minipage}[b]{\linewidth}\raggedright
origin
\end{minipage} \\
\midrule\noalign{}
\endhead
\bottomrule\noalign{}
\endlastfoot
\(\mathcal L_{\rm self}\) & --- & the self-interaction of the fermion
field (the author's definition, document KOHN\_SHAM\_GRAVITY):
\(\mathcal L_{\rm self}=-\tfrac12\sum_{r=0}^{8}g_r\,I_r\), a sum of nine
quartic invariants & K1 \\
\(g_r\); \(g_0,\dots,g_8\); \(g_S\), \(g_V\), \(g_A\), \(g_P\) &
\texttt{g\_0} \texttt{g\_1} \texttt{g\_2} \texttt{g\_3} \texttt{g\_4}
\texttt{g\_5} \texttt{g\_6} \texttt{g\_7} \texttt{g\_8} & the nine
couplings of the self-interaction, one for each rank \(r\) (the author's
symbols; dimension length\(^{D-1}\) in \(D\) space dimensions);
\(g_0=g_S\) (scalar), \(g_1=g_V\) (vector), \(g_7=-g_A\) (axial vector),
\(g_8=-g_P\) (pseudoscalar). Such a coupling carries one digit as
subscript; it is not a rate of deflation (\(\mathrm{g55}\): two digits,
upright, no subscript), not a component of the metric (superscript
\((4+4)\) or \((7+1)\), two indices) and not the coupling \(g_{VS}\) of
the Vorton-Scalar & K1 \\
\(K\), \(J\) (index sets); \(r\), \(r'\) (ranks) & \texttt{mask}
\texttt{rank} \texttt{J} & \(K\), \(J\): a strictly increasing list
\(A_1<\dots<A_r\) of indices out of \(0,\dots,7\) (there are 256);
\(r=\lvert K\rvert\) is its rank, \(0\le r\le8\); \(r'\) is a second
rank and \(l\) the summation index in the closed form of the Fierz
matrix. In the data files an index set is given by its bit mask
(\texttt{mask}: the sum of \(2^A\) over its indices) & K1 \\
\(\Theta^K\), \(\Theta_K\); \(\zeta_r\); \(\mathrm{sgn}_K\) &
\texttt{eps\_K} \texttt{phase} \texttt{eta\_K} & the 256 matrices of the
self-interaction (the author's symbols \(\Theta^K\), \(\Theta_K\)):
\(\Theta^K=\zeta_r\,\gamma^{A_1}\cdots\gamma^{A_r}\), with the phase
\(\zeta_r=1\) or \(i\) chosen so that the bilinear formed with
\(\Theta^K\) is real, and
\(\Theta_K=(\Theta^K)^{-1}=\mathrm{sgn}_K\,\Theta^K\) with a sign
\(\mathrm{sgn}_K=\pm1\). On 7+1 and in the quantized-sector form of 4+4:
\(\zeta_r=i\) for \(r=2,3,6,7\) (the author's \(i^{\sigma_r}\) with
\(\sigma_r=r(r-1)/2\) mod 2) and \(\mathrm{sgn}_K\) is the product of
the squares of the matrices in \(K\) (the author's \(\eta_K\)). In the
covariant form of 4+4: \(\zeta_r=i\) for \(r=1,2,5,6\), and
\(\mathrm{sgn}_K\) is \((-1)^r\) times that product & K1; the covariant
form of 4+4 is new \\
\(\gamma_{44q}^A\) & \texttt{gamma44q} & the matrices with which the
quantized-sector form of the self-interaction on 4+4 is built:
\(\gamma_{44q}^A=-i\,\gamma_{1644}^A\) for \(A=0,\dots,4\) and
\(\gamma_{44q}^A=\gamma_{1644}^A\) for \(A=5,6,7\). They have the
squares of the 7+1 matrices, and \(\gamma_{44q}^4=\beta_{1644}\) &
new \\
\(I_r\) & --- & the invariant of rank \(r\) (the author's symbol):
\(I_r=\sum_{\lvert K\rvert=r}(\bar\psi\Theta_K\psi)(\bar\psi\Theta^K\psi)\)
with \(\bar\psi=\psi^\dagger\beta\) on 7+1 and in the quantized-sector
form of 4+4, and with \(\tilde\psi=\psi^\dagger C^-_{1644}\) in place of
\(\bar\psi\) in the covariant form of 4+4. \(I_0\) is the square of the
scalar density, \(I_1\) the square of the current & K1 \\
\(\mathsf F\), \(\mathsf F_{rr'}\) & \texttt{fierz} \texttt{F\_r0}
\texttt{F\_r1} \texttt{F\_r2} \texttt{F\_r3} \texttt{F\_r4}
\texttt{F\_r5} \texttt{F\_r6} \texttt{F\_r7} \texttt{F\_r8} & the Fierz
matrix (\(9\times9\)):
\(\mathsf F_{rr'}=\tfrac1{16}\sum_l(-1)^{rr'-l}\binom{r'}{l}\binom{8-r'}{r-l}\),
with \(\mathsf F^2=1\). For anticommuting fields the nine invariants
obey \((1+\mathsf F)\,I=0\). (Sans-serif type; \(F_S\), \(F_P\),
\(F_0\), \(F_1\), \(F_7\), \(F_8\) in italic type are pairing
amplitudes. The author's document writes \(F_{rs}\).) & K1 \\
\(g_0^{\rm eff}\), \(g_1^{\rm eff}\), \(g_7^{\rm eff}\),
\(g_8^{\rm eff}\) & \texttt{g0eff} \texttt{g1eff} \texttt{g7eff}
\texttt{g8eff} \texttt{g\_0eff} \texttt{g\_1eff} \texttt{g\_7eff}
\texttt{g\_8eff} & the four independent couplings that remain for
anticommuting fields, when \(I_2,\dots,I_6\) are expressed by \(I_0\),
\(I_1\), \(I_7\), \(I_8\):
\(g_0^{\rm eff}=g_0-3g_2-4g_3-5g_4-4g_5-g_6\),
\(g_1^{\rm eff}=g_1-g_2-g_5+g_6\), \(g_7^{\rm eff}=g_7-g_2-g_3+g_6\),
\(g_8^{\rm eff}=g_8-g_2+4g_3-5g_4+4g_5-3g_6\); then
\(\mathcal L_{\rm self}=-\tfrac12\big(g_0^{\rm eff}I_0+g_1^{\rm eff}I_1+g_7^{\rm eff}I_7+g_8^{\rm eff}I_8\big)\)
& new \\
\(\mathsf n\) & --- & the one-body density matrix of matter
(\(16\times16\)), \(\mathsf n_{ab}=\langle\psi_b^\dagger\psi_a\rangle\),
with \(\mathrm{Tr}\,\mathsf n=n\) and
\(\mathrm{Tr}(\beta\,\mathsf n)=n_S\); for uniform matter at rest
\(\mathsf n=(n+n_S\,\beta)/16\). (The author's document writes \(R\) for
the matrix \(\mathsf n\,\beta\).) & new \\
\(G_S^{\rm HF}\), \(G_V^{\rm HF}\); \(G_S^{\rm c}\), \(G_V^{\rm c}\) &
\texttt{G\_0} \texttt{G\_1} \texttt{G\_S0} \texttt{G\_V0}
\texttt{hf\_rho} \texttt{hf\_rhoS} \texttt{hf\_rhoW} \texttt{hf\_rhoSW}
& the values that the Hartree-Fock approximation gives to the mean-field
couplings of uniform matter:
\(G_S^{\rm HF}=(15g_0^{\rm eff}-8g_1^{\rm eff}-8g_7^{\rm eff}-g_8^{\rm eff})/16\)
and
\(G_V^{\rm HF}=(-g_0^{\rm eff}+22g_1^{\rm eff}-6g_7^{\rm eff}+g_8^{\rm eff})/16\)
on 7+1 and in the quantized-sector form of 4+4, where the interaction
energy density is \(\tfrac12(G_V^{\rm HF}n^2+G_S^{\rm HF}n_S^2)\). In
the covariant form of 4+4 it is
\(\tfrac12(G_V^{\rm c}n^2+G_S^{\rm c}n_S^2+G_V^{\rm HF}n_W^2+G_S^{\rm HF}n_{SW}^2)\)
with \(G_V^{\rm c}=-(g_0^{\rm eff}+g_8^{\rm eff})/16\) and
\(G_S^{\rm c}=(-g_0^{\rm eff}-2g_1^{\rm eff}+2g_7^{\rm eff}+g_8^{\rm eff})/16\).
(The author's document and the data files write \(G_0\), \(G_1\) for
\(G_S^{\rm HF}\), \(G_V^{\rm HF}\); the data files write \texttt{G\_S0},
\texttt{G\_V0} for \(G_S^{\rm c}\), \(G_V^{\rm c}\).) & K1; the
covariant form of 4+4 is new \\
\(G_{\rm pair}\) & \texttt{G\_pair} & the contact coupling of a pair of
two particles at rest: the interaction energy of the pair times the
volume; attractive if negative. For the pair along the hidden axis
\(G_{\rm pair}=g_0^{\rm eff}+g_1^{\rm eff}+5\,g_7^{\rm eff}\) (7+1 and
quantized-sector form of 4+4) & new \\
\(n_+\), \(n_-\), \(n_{S+}\), \(n_{S-}\); \(n_W\), \(n_{SW}\);
\(\mathrm{pol}_W\) & \texttt{rho\_W} \texttt{rho\_SW} \texttt{pol\_W} &
number densities and scalar densities of the particles in the two halves
\(W_{1644}=\pm1\) of the 4+4 field; \(n_W=n_+-n_-\),
\(n_{SW}=n_{S+}-n_{S-}\); the polarization
\(\mathrm{pol}_W=(n_+-n_-)/n\) & new \\
\(\mathcal L_{\rm kin}\), \(\mathcal L_{\rm mass}\) & --- & the kinetic
term and the mass term of the Lagrangian of the fermion field,
\(\mathcal L=V_d\,[\mathcal L_{\rm kin}+\mathcal L_{\rm mass}+\mathcal L_{\rm self}]\)
& new \\
\(I_r^{\rm cov}\) & --- & the invariant of rank \(r\) of the covariant
form of the self-interaction on 4+4, written with the matrices of the
quantized-sector form:
\(I_r^{\rm cov}=\sum_{\lvert K\rvert=r}\mathrm{sgn}_K(\psi^\dagger W_{1644}\beta\Theta^K\psi)^2\)
& new \\
\end{longtable}

\textbf{12 Kohn-Sham and CJT calculations with contact interaction}

\begin{longtable}[]{@{}
  >{\raggedright\arraybackslash}p{(\columnwidth - 6\tabcolsep) * \real{0.1605}}
  >{\raggedright\arraybackslash}p{(\columnwidth - 6\tabcolsep) * \real{0.1235}}
  >{\raggedright\arraybackslash}p{(\columnwidth - 6\tabcolsep) * \real{0.5309}}
  >{\raggedright\arraybackslash}p{(\columnwidth - 6\tabcolsep) * \real{0.1852}}@{}}
\toprule\noalign{}
\begin{minipage}[b]{\linewidth}\raggedright
symbol
\end{minipage} & \begin{minipage}[b]{\linewidth}\raggedright
names in programs and data files
\end{minipage} & \begin{minipage}[b]{\linewidth}\raggedright
meaning
\end{minipage} & \begin{minipage}[b]{\linewidth}\raggedright
origin
\end{minipage} \\
\midrule\noalign{}
\endhead
\bottomrule\noalign{}
\endlastfoot
\(G_S\), \(G_V\) & \texttt{G\_S} \texttt{G\_V} \texttt{GS\_abs} & the
scalar and the vector coupling of the Kohn-Sham (mean-field) description
of uniform matter (the author's symbols): interaction energy density
\(\tfrac12G_Vn^2+\tfrac12G_Sn_S^2\); attraction for \(G_S<0\), repulsion
for \(G_V>0\); dimension length\(^{D-1}\). The author's example:
\(G_S=-10^5\), \(G_V=+10^4\) in units \(m=1\), \(D=7\) & K1 \\
\(n\), \(n_S\) & \texttt{rho} \texttt{rho\_S} \texttt{rho\_ks}
\texttt{rho\_cjt} & number density \(\langle\psi^\dagger\psi\rangle\)
and scalar density \(\langle\bar\psi\psi\rangle\) of the fermions of
matter (the filled states of negative energy not counted). The data
files and the author's document write \(\rho\) (\texttt{rho}) and
\(\rho_S\) (\texttt{rho\_S}) & K1 (there \(\rho\), \(\rho_S\)) \\
\(D\) & \texttt{D} & the number of space dimensions in which the
particles of matter have momentum: 7 (all space coordinates of 7+1), 4
(fields that do not depend on \(x5,x6,x7\)) or 3 (fields that depend on
ordinary space only). \(V_D=\pi^{D/2}/\Gamma_{\rm E}(D/2+1)\) is the
volume and \(A_D=D\,V_D\) the area of the unit ball in \(D\) dimensions
(\(\Gamma_{\rm E}\): Euler's gamma function) & new \\
\(E^*(k)\); \(\langle f\rangle\) & --- & \(E^*(k)=\sqrt{k^2+m^{*2}}\),
the energy of a particle of matter without the vector term;
\(\langle f\rangle=\nu\int d^Dk\,f/(2\pi)^D\) over the occupied states:
\(n=\langle1\rangle\), \(n_S=\langle m^*/E^*\rangle\) & K1 (there
\(\omega\) for \(E^*\)) \\
\(\varepsilon_{\rm kin}\); \(p_j\); \(p_{\rm int}\) & \texttt{e\_kin}
\texttt{eps} \texttt{eps\_ks} \texttt{eps\_cjt} \texttt{eps\_rel}
\texttt{p} \texttt{p\_ks} \texttt{p\_flat} \texttt{p5} \texttt{p6}
\texttt{p7} \texttt{p\_a} \texttt{p\_b} \texttt{pa\_rel}
\texttt{pb\_rel} \texttt{Pi} &
\(\varepsilon_{\rm kin}=\langle E^*\rangle\): the kinetic part of the
energy density
\(\varepsilon=\varepsilon_{\rm kin}+\tfrac12G_Vn^2-(m^*-m)^2/(2G_S)\);
\(p_j\): the pressure of matter along the coordinate \(xj\),
\(\langle k_j^2/E^*\rangle+p_{\rm int}\);
\(p_{\rm int}=\tfrac12G_Vn^2+(m^*-m)^2/(2G_S)\): the interaction part of
the pressure, which is the whole pressure along a coordinate on which
the fields do not depend (\texttt{Pi} in the data files, \(\Pi\) in the
author's document) & K1 \\
\(V_0\) & \texttt{V\_0} & the vector mean field \(V_0=G_V\,n\) (the
author's symbol); it shifts the chemical potential, \(\mu\to\mu-V_0\) &
K1 \\
\(\Omega_{\rm CJT}\) & \texttt{minus\_Omega} & the CJT effective
potential of uniform matter with contact interaction per unit volume, a
function of the mass \(m^*\) and of the vector field \(V_0\) at the
chemical potential \(\mu\):
\(\Omega_{\rm CJT}=-\nu\int\frac{d^Dk}{(2\pi)^D}\,(\mu-V_0-E^*)\,\theta(\mu-V_0-E^*)-\frac{(m^*-m)^2}{2G_S}-\frac{V_0^2}{2G_V}\);
at its stationary point \(-\Omega_{\rm CJT}\) is the pressure & new \\
\(G_c\) & \texttt{Gc} & the critical coupling of the vacuum: with the
filled states of negative energy below the cutoff \(\Lambda\) a massless
field acquires a mass when
\(G_S<G_c=-(2\pi)^D(D-1)/(\nu A_D\Lambda^{D-1})\);
\(G_c\Lambda^6=-90\pi^4\) for \(D=7\) & K2 (\(D=7\)) \\
\(m_{\rm vac}\), \(m_{\rm bare}\);
\(\langle\bar\psi\psi\rangle_{\rm vac}\) & \texttt{M\_vac}
\texttt{m\_bare} \texttt{condensate} & with the filled states of
negative energy: the mass of the fermions in vacuum, the mass parameter
of the Lagrangian (bare mass) that produces it, and the scalar density
of the filled states,
\(\langle\bar\psi\psi\rangle_{\rm vac}=-\nu\int^\Lambda\frac{d^Dk}{(2\pi)^D}\frac{m_{\rm vac}}{\sqrt{k^2+m_{\rm vac}^2}}\)
(negative);
\(m_{\rm vac}=m_{\rm bare}+G_S\langle\bar\psi\psi\rangle_{\rm vac}\) &
new \\
\(m_D\), \(m_Y\), \(f_D\) & \texttt{f\_D} & when a bare Dirac mass and
the Vorton-Scalar are both present: the bare Dirac mass \(m_D\), the
part \(m_Y=g_{VS}v_{VS}\) of the mass that the Vorton-Scalar gives, and
the fraction \(f_D=m_D/m\) of the vacuum mass \(m=m_D+m_Y\) & new \\
\(m_\vartheta\) & \texttt{mph2\_over\_C\_ms2} & the mass of the phase
mode \(\vartheta_{VS}\) of the Vorton-Scalar: zero while the chiral
symmetry \(U(1)_{VS}\) is exact; a bare Dirac mass breaks the symmetry
and gives
\(m_\vartheta^2v_{VS}^2=\big(\lvert\langle\bar\psi\psi\rangle_{\rm vac}\rvert-n_S\big)m_Dm_Y/m\)
& new \\
\(-G_S/C\), \(G_V/C\) & \texttt{kappa\_S} \texttt{kappa\_V} & the
contact couplings in units of the coupling \(C=g_{VS}^2m^2/m_s^2\) of
the Vorton-Scalar (units \(m=1\)); they are written as these ratios,
without a letter of their own (columns \texttt{kappa\_S},
\texttt{kappa\_V}) & new \\
\(\mathsf h_{\rm KS}\) & --- & the Hamiltonian matrix of a particle of
matter in the Kohn-Sham description,
\(\mathsf h_{\rm KS}=\sum_j\alpha^jk_j+\beta\,m^*+V_0\) & K1 \\
\(\Gamma[S]\); \(S\), \(S_0\) & --- & the CJT effective action as a
functional of the propagator of the fermions,
\(\Gamma[S]=-i\,\mathrm{Tr}\ln S^{-1}-i\,\mathrm{Tr}(S_0^{-1}S-1)+\Gamma_2[S]\),
with the full propagator \(S\) and the propagator \(S_0\) of free
particles of the mass \(m\) (the argument in brackets tells it from the
action \(\Gamma[\phi,G]\) of the model with the Vorton-Scalar) & new \\
\(n_{\rm ref}\) & --- & the number density of matter at the reference
time \(x4_{\rm ref}\); matter that keeps its particles has
\(n=n_{\rm ref}\,e^{\tau_d}\) & new \\
\end{longtable}

\textbf{13 Pairing with contact couplings}

\begin{longtable}[]{@{}
  >{\raggedright\arraybackslash}p{(\columnwidth - 6\tabcolsep) * \real{0.1605}}
  >{\raggedright\arraybackslash}p{(\columnwidth - 6\tabcolsep) * \real{0.1235}}
  >{\raggedright\arraybackslash}p{(\columnwidth - 6\tabcolsep) * \real{0.5309}}
  >{\raggedright\arraybackslash}p{(\columnwidth - 6\tabcolsep) * \real{0.1852}}@{}}
\toprule\noalign{}
\begin{minipage}[b]{\linewidth}\raggedright
symbol
\end{minipage} & \begin{minipage}[b]{\linewidth}\raggedright
names in programs and data files
\end{minipage} & \begin{minipage}[b]{\linewidth}\raggedright
meaning
\end{minipage} & \begin{minipage}[b]{\linewidth}\raggedright
origin
\end{minipage} \\
\midrule\noalign{}
\endhead
\bottomrule\noalign{}
\endlastfoot
\(F_0\), \(F_1\), \(F_7\), \(F_8\) & \texttt{F0} \texttt{F1} \texttt{F7}
\texttt{F8} \texttt{F0c} \texttt{F1c} \texttt{F7c} \texttt{F8c} & the
pairing amplitudes of the four independent contact channels in the pair
channel along the hidden axis, for particles of mass \(m^*\) that
scatter from the momenta \(\vec p,-\vec p\) to \(\vec p',-\vec p'\) in
ordinary space (\(E=\sqrt{m^{*2}+p^2}\), \(E'\) likewise):
\(F_0=(EE'+m^{*2}-\vec p\cdot\vec p')/(2EE')\),
\(F_1=(3EE'-2m^{*2}-\vec p\cdot\vec p')/(EE')\),
\(F_7=(3EE'+2m^{*2}-\vec p\cdot\vec p')/(EE')\),
\(F_8=-(EE'-m^{*2}-\vec p\cdot\vec p')/(2EE')\). \(F_0=F_S\) and
\(F_8=-F_P\) & new \\
\(m_{\rm ct}\) & \texttt{m\_ct} & the range parameter of the contact
couplings in the gap equation: each contact coupling \(g\) is replaced
by \(g\,m_{\rm ct}^2/(q^2+m_{\rm ct}^2)\), the form that the exchange of
a boson of mass \(m_{\rm ct}\) gives (\(q\): momentum transfer) & new \\
\(V_{\rm ct}\), \(V_{VS}\) & --- & the pairing potentials in the
hidden-axis channel: of the contact interaction,
\(V_{\rm ct}=\sum_{r=0,1,7,8}g_r^{\rm eff}F_r\,m_{\rm ct}^2/(q^2+m_{\rm ct}^2)\),
and of the exchange of the two modes of the Vorton-Scalar,
\(V_{VS}=-g_{VS}^2\big[F_S/(q^2+m_s^2)+F_P/(q^2+m_\vartheta^2)\big]\);
attractive where negative & new \\
\(\Phi\), \(\Phi(\vec p)\) & --- & a state of a pair of two particles,
an antisymmetric \(16\times16\) matrix; \(\Phi(\vec p)\): the state of
the pair along the hidden axis with the momenta \(\vec p\), \(-\vec p\),
the matrix \(X^*=(C^-\gamma^0)^*\) between the projectors on the states
of positive energy. (Without suffix; the Vorton-Scalar always carries
the suffix VS.) & new \\
\end{longtable}

\textbf{14 Einstein-Lovelock gravity with Kohn-Sham matter}

\begin{longtable}[]{@{}
  >{\raggedright\arraybackslash}p{(\columnwidth - 6\tabcolsep) * \real{0.1605}}
  >{\raggedright\arraybackslash}p{(\columnwidth - 6\tabcolsep) * \real{0.1235}}
  >{\raggedright\arraybackslash}p{(\columnwidth - 6\tabcolsep) * \real{0.5309}}
  >{\raggedright\arraybackslash}p{(\columnwidth - 6\tabcolsep) * \real{0.1852}}@{}}
\toprule\noalign{}
\begin{minipage}[b]{\linewidth}\raggedright
symbol
\end{minipage} & \begin{minipage}[b]{\linewidth}\raggedright
names in programs and data files
\end{minipage} & \begin{minipage}[b]{\linewidth}\raggedright
meaning
\end{minipage} & \begin{minipage}[b]{\linewidth}\raggedright
origin
\end{minipage} \\
\midrule\noalign{}
\endhead
\bottomrule\noalign{}
\endlastfoot
\(E^{(k)\mu}{}_\nu\) & --- & the Lovelock tensor of order \(k=1,2,3\) in
the author's notation and sign: \(E^{(1)}=-\mathcal G\) (minus the
Einstein tensor), \(E^{(2)}=\mathcal H\) (the Lanczos tensor). For a
metric whose scale factors depend only on \(x4\):
\(E^{(k)m}{}_m=(-1)^{k+1}k!\,2^{k-1}\) times the sum, over the sets of
\(2k\) frame directions other than \(m\) and over their divisions into
pairs, of the products of the sectional curvatures \(K_{AB}\) of the
pairs. (With the order in parentheses as a superscript; \(E_F\),
\(E_k\), \(E_\perp\), \(E^*\) are energies.) & K1 \\
\(\alpha_1\), \(\alpha_2\), \(\alpha_3\); \(\bar\Lambda\); \(\kappa_8\)
& \texttt{LambdaBar} \texttt{kappa} \texttt{kappa\_eps} & the couplings
of the Lovelock terms (\(\alpha_1=1\)), the cosmological term and the
gravitational coupling of the author's Einstein-Lovelock equations:
\(\sum_k\alpha_kE^{(k)4}{}_4-\bar\Lambda=\kappa_8\,\varepsilon\) for the
time component and
\(\sum_k\alpha_kE^{(k)m}{}_m-\bar\Lambda=-\kappa_8\,p_m\) for the space
direction \(m\). \(\alpha_2=-\alpha_{GB}\) and
\(\bar\Lambda=\Lambda_8\). The author's document writes \(\kappa\) for
\(\kappa_8=8\pi G_8\) (the letter \(\kappa\) with an argument,
\(\kappa(k)\), is a pair correlation). A Lovelock coupling carries a
digit as subscript; \(\alpha^j\) with a superscript are Dirac matrices &
K1 \\
\(\mathsf a\), \(\mathsf b\); \(H_{\mathsf a}\), \(H_{\mathsf b}\);
\(n_{\mathsf a}\), \(n_{\mathsf b}\) & \texttt{a} \texttt{b}
\texttt{H\_a} \texttt{H\_b} \texttt{Ha\_over\_g} \texttt{Hb\_over\_g}
\texttt{Hb0\_over\_g} \texttt{a\_end} \texttt{b\_end}
\texttt{Ha\_end\_over\_g} \texttt{Hb\_end\_over\_g} & the two scale
factors of a spacetime whose space coordinates form two groups of
\(n_{\mathsf a}\) and \(n_{\mathsf b}\) coordinates, and their rates
\(H_{\mathsf a}=\mathsf a'/\mathsf a\),
\(H_{\mathsf b}=\mathsf b'/\mathsf b\) (a prime: derivative with respect
to the time \(x4\)). In the author's document: \(\mathsf a\) for
ordinary space (\(n_{\mathsf a}=3\)) and \(\mathsf b\) for the four
auxiliary coordinates. On the deflating spacetimes with equal rates:
\(\mathsf a\) for \(x0,\dots,x3\) (\(n_{\mathsf a}=4\)) and
\(\mathsf b=b_5=b_6=b_7\) (\(n_{\mathsf b}=3\)); the undisturbed
deflation is \(H_{\mathsf a}=0\), \(H_{\mathsf b}=-\mathrm{g55}\) & K1
(there \(a\), \(b\), \(H_a\), \(H_b\)) \\
loading & \texttt{loading} &
\(\kappa_8\,\varepsilon(0)/(3\,\mathrm{g55}^2)\): the energy density of
the matter at the start of an evolution, in units of the energy density
that would make three coordinates deflate at the rate \(\mathrm{g55}\)
in Einstein gravity & new \\
\(T^\mu{}_\nu\); \(p_{\mathsf a}\), \(p_{\mathsf b}\);
\(\delta H_{\mathsf b}\) & --- & the energy-momentum tensor of uniform
matter in mixed components, \(T^4{}_4=\varepsilon\), \(T^m{}_m=-p_m\);
its pressures along the coordinates that carry the scale factors
\(\mathsf a\) and \(\mathsf b\); the deviation
\(H_{\mathsf b}+\mathrm{g55}\) of the rate of the deflating coordinates
from the undisturbed deflation & new \\
\end{longtable}

\textbf{15 Tags and identifiers of the calculations with mass term and
self-interaction}

\begin{longtable}[]{@{}
  >{\raggedright\arraybackslash}p{(\columnwidth - 6\tabcolsep) * \real{0.1605}}
  >{\raggedright\arraybackslash}p{(\columnwidth - 6\tabcolsep) * \real{0.1235}}
  >{\raggedright\arraybackslash}p{(\columnwidth - 6\tabcolsep) * \real{0.5309}}
  >{\raggedright\arraybackslash}p{(\columnwidth - 6\tabcolsep) * \real{0.1852}}@{}}
\toprule\noalign{}
\begin{minipage}[b]{\linewidth}\raggedright
symbol
\end{minipage} & \begin{minipage}[b]{\linewidth}\raggedright
names in programs and data files
\end{minipage} & \begin{minipage}[b]{\linewidth}\raggedright
meaning
\end{minipage} & \begin{minipage}[b]{\linewidth}\raggedright
origin
\end{minipage} \\
\midrule\noalign{}
\endhead
\bottomrule\noalign{}
\endlastfoot
71s, 44q, 44c & \texttt{71s} \texttt{44q} \texttt{44c} & tags of the
three models of the self-interaction: 71s on 7+1; 44q the
quantized-sector form on 4+4 (bilinears with
\(\bar\psi=\psi^\dagger\beta_{1644}\) and the matrices
\(\gamma_{44q}^A\)); 44c the covariant form on 4+4 (bilinears with
\(\tilde\psi=\psi^\dagger C^-_{1644}\)) & new \\
identifiers of the machine checks & --- & \texttt{44d-..},
\texttt{71d-..}, \texttt{def-..}: checks on the deflating 4+4 spacetime,
on the deflating 7+1 spacetime, and of results that hold on both;
\texttt{44-..}, \texttt{71-..}: checks of the flat spacetimes, executed
again; \texttt{71s-S..}, \texttt{44q-S..}, \texttt{44c-S..},
\texttt{self-X..}: the algebra of the self-interaction; \texttt{ks-K..}:
Kohn-Sham calculations; \texttt{cjt-J..}: CJT effective potential with
contact interaction; \texttt{uni-U..}: contact interaction, bare mass
and Vorton-Scalar together; \texttt{pc-P..}: pairing with contact
couplings; \texttt{gr-G..}: gravity with Kohn-Sham matter. A leading
capital I (\texttt{Icks-..}, \texttt{Idef-..}, \texttt{I44d-..},
\texttt{I71d-..}) marks a check of the independent verification in
Python & new \\
\end{longtable}

\hypertarget{key-to-the-column-origin}{%
\section{8. Key to the column
``origin''}\label{key-to-the-column-origin}}

\begin{longtable}[]{@{}
  >{\raggedright\arraybackslash}p{(\columnwidth - 2\tabcolsep) * \real{0.5000}}
  >{\raggedright\arraybackslash}p{(\columnwidth - 2\tabcolsep) * \real{0.5000}}@{}}
\toprule\noalign{}
\begin{minipage}[b]{\linewidth}\raggedright
code
\end{minipage} & \begin{minipage}[b]{\linewidth}\raggedright
source
\end{minipage} \\
\midrule\noalign{}
\endhead
\bottomrule\noalign{}
\endlastfoot
N01 & the author's notebook \texttt{4+4.nb} (Dirac matrices of 4+4;
split octonions) \\
N02 & the author's notebook \texttt{7+1.nb} (octonions; Dirac matrices
and Dirac equation of 7+1) \\
N04 & the author's notebook
\texttt{Dirac\_Einstein-Lovelock\_4+4\_Nash.nb} (meaning of the
coordinates) \\
N05 & the author's notebook \texttt{Dirac\_Nash\_Identity.nb} \\
N10 & the author's notebook
\texttt{Pair\_Creation\_of\_Universes\_WaveFunctionOfUniverse-4+4-Einstein-Lovelock-Nash.nb} \\
N12 & the author's notebook
\texttt{WaveFunctionUniverse-4+4-Einstein-Lovelock.nb} \\
D06 & P. L. Nash, ``Identities satisfied by the generators of the Dirac
algebra'', J. Math. Phys. 25, 204 (1984) \\
D07 & P. L. Nash, ``On the structure of the split octonion algebra'',
Nuovo Cimento B 105, 31 (1990) \\
D08 & P. L. Nash, ``Second gravity'', J. Math. Phys. 51, 042501
(2010) \\
K1 & the author's document \texttt{KOHN\_SHAM\_GRAVITY} (folder 7+1;
sections 15 and 16 of the notebook \texttt{7+1.nb}): nine-channel
self-coupling, Fierz matrix, Kohn-Sham theory, Einstein-Lovelock
evolution \\
K2 & the author's document \texttt{DIRAC16COMPLEX} (folder 7+1): the
critical coupling of the vacuum in seven space dimensions \\
author & a name that an instruction of the author fixes (the rates
\(\mathrm{g55}\), \(\mathrm{g66}\), \(\mathrm{g77}\); the constants; the
signs of the metric) \\
new & introduced by this project \\
\end{longtable}

\end{document}
